% ==============================================================================
% METODOLOGÍA
% ==============================================================================
% Evolución del Marco Metodológico de protocolo_sebastian_v1 (3).pdf
% (sección 7), actualizada al estado real del proyecto al 2026-10-02.
%
% Fuentes de cada cifra y procedimiento (repositorio logSguarDian y
% logSguarDian-vulnerable-project):
%   - Dataset y partición: training/split.py, training/DEDUP_METHODOLOGY.md,
%     docs/dataset-audit.md, training/results/v11_test_results.json.
%   - Umbral RF: docs/decision-policy.md secciones 2.2.1 y 2.2.2.
%   - Corpus limpio (Ronda 5): attack-sim/build_corpus.py,
%     attack-sim/exclusion_filter.py, docs/config3b-results.md.
%   - Fuga: training/leakage_check.py, training/results/leakage_report.json.
%   - Latencia/memoria/CPU: .github/workflows/latency-benchmark.yml,
%     attack-sim/run_latency_suite.sh, artillery-*.yml, monitor-memory.sh.
%   - Apps de terceros: docs/vulnerable-projects-tests/,
%     docs/findings-evidence/juice-shop-ua-bias/.
%
% Decisiones del estudiante incorporadas (conversación 2026-10-02):
%   - Criterio de latencia absoluto de 15 ms (ya fijado en h-alcance.tex),
%     justificado porque una cifra en milisegundos es más representativa
%     que un porcentaje sobre una línea base de pocos milisegundos.
%   - Umbral RF calibrado con rf_v3.
%   - Tablero Kanban de cuatro columnas.
%   - Kanban en Notion + flujo main <- develop <- rama con aprobación
%     restringida a los dos integrantes + GitHub Actions.
%   - La validación con startups no se realizó; se sustituye por la
%     validación en aplicaciones de terceros.
%   - [Actualizado 2026-10-10: ZAP/XSS y Config 1/3a/3b ya se ejecutaron sobre el corpus limpio.]
%     Antes pendientes: (ZAP/XSS, Config 1/3a/3b con corpus
%     limpio, latencia con el modelo sincronizado, memoria en escenario
%     único) no se ejecutarán; se declaran como restricciones.
%   - El paquete mlops se incorpora por recomendación del asesor, del lado
%     del componente de IA; aquí solo se describe como interfaz.
%
% Fechas de la Tabla de fases: reconstruidas del historial de git de ambos
% repositorios (primer y último commit relacionado). Pendientes de
% confirmación por el estudiante.
\ifdefined\modoensamble
\else
	\documentclass[11pt, letterpaper]{report}
\usepackage[spanish, es-nodecimaldot, es-noquoting]{babel}
\selectlanguage{spanish}
\usepackage[utf8]{inputenc}
\usepackage[T1]{fontenc}

% ===============================================================================
% El estudiante debe llenar sus datos en esta sección para que la plantilla los 
% auto-importe y genere automáticamente las páginas de portada y de firmas 
% autorizadas.
% ===============================================================================
% Datos del estudiante:
% -------------------------------------------------------------------------------
% Nombre completo
\def \nombreestudiante {Sebastián Huertas Gómez}
% Carné
\def \uvgcarne {22295}
% Facultad
\def \uvgfacultad {Ingeniería}
% Carrera
\def \uvgcarrera {Ingeniería en Ciencia de la Computación y Tecnologías de la Información}

% Datos del trabajo:
% -------------------------------------------------------------------------------
% Título completo (tomado literalmente de protocolo_sebastian_v1 (3).pdf, sección 1)
\def \titulotesis {logSguarDian: Desarrollo de una librería para la detección de amenazas de seguridad en Endpoints de aplicaciones}
% Mes y año de entrega (confirmado por el estudiante: noviembre 2026)
\def \mesentrega {noviembre}
\def \anoentrega {2026}
% Asesor (según protocolo_sebastian_v1 (3).pdf, portada)
\def \nombreasesor {Ing. Erick Francisco Marroquín Rodríguez}
% El catedrático del curso (Ing. Bacilio Alexander Bolaños Lima) figura como
% primer examinador en la hoja de firmas, definido más abajo.

% Datos del tribunal examinador:
% -------------------------------------------------------------------------------
% Nombre del primer examinador (segunda línea de la hoja de firmas)
\def \nombreprimerex {Ing. Ma. Bacilio Alexander Bolaños Lima}

% Capítulos pre-definidos
% -------------------------------------------------------------------------------
% Comentar las líneas de las secciones que desean omitirse, por defecto se
% se incluyen todas.
% Se descomentan conforme se redacta cada capítulo (todavía sin archivo de
% contenido: prefacio, alcance, anexos, glosario).
\def \CAPprefacio {Prefacio}
\def \CAPantecedentes {Antecedentes}
\def \CAPalcance {Alcance}
\def \CAPmetodologia {Metodología}
\def \CAPresultados {Resultados}
\def \CAPdiscusion {Análisis de resultados}
\def \CAPanexos {Anexos}
\def \CAPglosario {Glosario}

% Formato y estilo de la plantilla
% -------------------------------------------------------------------------------
% Modo impresión: Puede des-comentar la siguiente línea para generar un documento pdf sin la portada, para cuando se desee imprimir el documento para encuadernación
% NOTA: la plantilla original traía esta línea activa por defecto, lo cual
% suprime la portada verde (main.tex hace \let\CAPportada\undefined cuando
% \printver está definido). Se comenta para que la versión digital, con
% portada, sea la que se genera por defecto.
%\def \printver {Versión del documento para impresión}

% Portada: Puede cambiarse la imagen en la portada al cambiar el nombre del 
% archivo siguiente. NOTA: debe tener la suficiente resolución para cubrir el área
% designada
\def \imagenportada {plantilla/portadacit.jpg}

% Referencias: Puede des-comentar la siguiente línea para utilizar el formato de referencias APA
%\def \usarAPA {Usar formato APA}

% Párrafo: Puede comentar la siguiente línea si desea emplear un formato de 
% párrafo distinto al establecido por defecto
\def \parpordefecto {Formato de párrafo por defecto}

% Capítulos y secciones: Puede des-comentar la siguiente línea para establecer el
% formato de los capítulos y secciones bajo el estándar original de UVG para
% trabajos de graduación. Este incluye: capítulos con numeración romana, secciones
% con letras mayúsculas, sub-secciones con números y sub-sub-secciones con letras
% minúsculas
%\def \capsecuvg {Formato UVG para capítulos y secciones}
% ================================================================================
% En este archivo se colocan opciones adicionales para modificar el formato de la
% plantilla, para emplearse en otros tipos de documentos que no sean trabajos de
% graduación. Si usted está trabajando su tesis, NO modifique este archivo
% ================================================================================
% Capítulos pre-definidos
% --------------------------------------------------------------------------------
% Comentar las líneas de las secciones que desean omitirse, por defecto se 
% se incluyen todas.
\def \CAPportada {Portada}
\def \CAPcaratula {Caratula}
\def \CAPfirmas {Hoja de firmas}
\def \CAPindice {Índice general}
\def \CAPfiguras {Listado de figuras}
\def \CAPcuadros {Listado de cuadros}
\def \CAPresumen {Resumen}
\def \CAPabstract {Resumen}
\def \CAPintroduccion {Introducción}
\def \CAPobjetivos {Objetivos}
\def \CAPjustificacion {Justificación}
\def \CAPmarcoteorico {Marco teórico}
\def \CAPconclusiones {Conclusiones}
\def \CAPrecomendaciones {Recomendaciones}
\def \CAPbibliografia {Bibliografía}

\usepackage{xcolor}
\usepackage{amsfonts}
\usepackage{amsmath}
\usepackage{amssymb}
\usepackage{amsthm}
\usepackage{mathtools}
\usepackage{graphicx}
\usepackage{xfrac}
\usepackage{float}
\usepackage[hypertexnames=false]{hyperref}
\usepackage[font=small]{caption}
\usepackage{subcaption}
% Cuadros: array (columnas L{} alineadas a la izquierda) y longtable (cuadros
% que ocupan más de una página). Se cargan aquí, no en
% 2-paquetes_y_comandos_usuario.tex, para que los capítulos de Metodología y
% Resultados no dependan de ese archivo.
\usepackage{array}
\usepackage{longtable}
\newcolumntype{L}[1]{>{\raggedright\arraybackslash}p{#1}}
\renewcommand{\arraystretch}{1.3}
\setlength{\LTpre}{0.3cm}
\setlength{\LTpost}{0.1cm}
% Margen de estiramiento para que los párrafos con palabras largas no
% generen líneas que sobresalgan del margen.
\setlength{\emergencystretch}{3em}
\usepackage{xpatch}
\usepackage{emptypage}
% htt: permite dividir con guion el texto en \texttt (nombres de archivo,
% clases y campos), que de otro modo sobresale del margen derecho.
\usepackage[htt]{hyphenat}
\usepackage{fancyhdr}
\usepackage[backend=biber, style=ieee]{biblatex}
\ifdefined\usarAPA
    \usepackage[backend=biber, style=apa]{biblatex}
\fi
\addbibresource{m-bibliografia.bib}

\usepackage[percent]{overpic}
\usepackage{chngcntr}

\ifdefined\CAPglosario
	\usepackage[numberedsection]{glossaries}
	\makeglossaries
    \newglossaryentry{rasp}
{
    name={RASP},
    description={Autoprotección de aplicaciones en tiempo de ejecución (\textit{Runtime Application Self-Protection}): mecanismo de seguridad que se ejecuta dentro de la propia aplicación y analiza sus peticiones para detectar y bloquear ataques.}
}

\newglossaryentry{waf}
{
    name={WAF},
    description={Cortafuegos de aplicaciones web (\textit{Web Application Firewall}): dispositivo o servicio perimetral que filtra el tráfico HTTP antes de que llegue a la aplicación.}
}

\newglossaryentry{modsecurity}
{
    name={ModSecurity},
    description={Motor de WAF de código abierto que se usa con reglas externas; en este trabajo se configuró con OWASP CRS.}
}

\newglossaryentry{crs}
{
    name={OWASP CRS},
    description={Conjunto de reglas genéricas de OWASP (\textit{Core Rule Set}) para ModSecurity, con niveles de paranoia que controlan qué tan estrictas son.}
}

\newglossaryentry{owasp}
{
    name={OWASP Top 10},
    description={Lista de las diez categorías de riesgo más críticas para aplicaciones web, publicada por OWASP; este trabajo usa la edición 2021.}
}

\newglossaryentry{mitre}
{
    name={MITRE ATT\&CK},
    sort={MITRE ATT&CK},
    description={Base de conocimiento de tácticas y técnicas de atacantes, usada aquí para asociar cada vector de ataque a una técnica.}
}

\newglossaryentry{cwe}
{
    name={CWE},
    description={Enumeración común de debilidades de software (\textit{Common Weakness Enumeration}): catálogo de tipos de defectos que originan vulnerabilidades.}
}

\newglossaryentry{sqli}
{
    name={SQL Injection},
    description={Ataque que inserta sentencias SQL en una entrada de la aplicación para alterar las consultas que ejecuta la base de datos.}
}

\newglossaryentry{xss}
{
    name={XSS},
    description={\textit{Cross-Site Scripting}: ataque que inyecta código ejecutable en una página que otros usuarios verán.}
}

\newglossaryentry{pathtraversal}
{
    name={Path Traversal},
    description={Ataque que manipula rutas de archivo (por ejemplo con \texttt{../}) para leer archivos fuera del directorio previsto; incluye la inclusión local de archivos (LFI).}
}

\newglossaryentry{cmdi}
{
    name={Command Injection},
    description={Ataque que logra ejecutar comandos del sistema operativo a través de una entrada de la aplicación.}
}

\newglossaryentry{middleware}
{
    name={\textit{middleware}},
    sort={middleware},
    description={Función que se ejecuta entre la recepción de una petición y la lógica de negocio de la aplicación; en Express se registra con \texttt{app.use}.}
}

\newglossaryentry{rf}
{
    name={Random Forest},
    description={Algoritmo supervisado que combina muchos árboles de decisión; en la librería es el único modelo con autoridad para bloquear una petición.}
}

\newglossaryentry{iforest}
{
    name={Isolation Forest},
    description={Algoritmo no supervisado de detección de anomalías que aísla las observaciones atípicas con particiones aleatorias; en la librería solo registra anomalías y nunca bloquea.}
}

\newglossaryentry{onnx}
{
    name={ONNX},
    description={Formato abierto para serializar modelos de aprendizaje automático, que permite ejecutarlos fuera de Python.}
}

\newglossaryentry{workerthreads}
{
    name={\texttt{worker\_threads}},
    sort={worker_threads},
    description={Módulo de Node.js que ejecuta código en hilos independientes del hilo principal; la inferencia de los modelos se ejecuta ahí.}
}

\newglossaryentry{eventloop}
{
    name={\textit{Event Loop}},
    sort={Event Loop},
    description={Bucle de eventos de Node.js que atiende las peticiones en un único hilo; bloquearlo degrada a toda la aplicación.}
}

\newglossaryentry{failopen}
{
    name={\textit{fail-open}},
    sort={fail-open},
    description={Política por la cual, ante un error o un tiempo de espera del módulo de detección, la petición se deja pasar en lugar de bloquearse.}
}

\newglossaryentry{monitor}
{
    name={Modo \texttt{monitor}},
    sort={Modo monitor},
    description={Modo de operación en el que la librería registra las peticiones que bloquearía, pero no las bloquea.}
}

\newglossaryentry{block}
{
    name={Modo \texttt{block}},
    sort={Modo block},
    description={Modo de operación en el que la librería responde HTTP 403 a las peticiones que clasifica como ataque.}
}

\newglossaryentry{pesp}
{
    name={p95},
    description={Percentil 95 del tiempo de respuesta: el valor por debajo del cual queda el 95~\% de las peticiones.}
}

\newglossaryentry{deltap95}
{
    name={$\Delta p95$},
    sort={Delta p95},
    description={Diferencia entre el p95 con el middleware activo y el p95 de la línea base sin él; en este trabajo se mide en milisegundos.}
}

\newglossaryentry{funo}
{
    name={F1},
    description={Media armónica de la precisión y el recall; es la métrica principal de clasificación del trabajo.}
}

\newglossaryentry{auc}
{
    name={AUC-ROC},
    description={Área bajo la curva ROC; mide qué tan bien separa un modelo dos clases sin depender de un umbral.}
}

\newglossaryentry{wilson}
{
    name={Intervalo de Wilson},
    description={Intervalo de confianza para una proporción, adecuado para muestras pequeñas o proporciones cercanas a 0 o 1.}
}

\newglossaryentry{casidup}
{
    name={Casi-duplicado},
    description={Registro muy parecido a otro pero no idéntico; si queda repartido entre entrenamiento y prueba, puede inflar las métricas.}
}

\newglossaryentry{minhash}
{
    name={MinHash},
    description={Técnica para estimar rápidamente la similitud entre conjuntos, usada para detectar casi-duplicados.}
}

\newglossaryentry{fuga}
{
    name={Fuga de datos},
    description={Presencia en el conjunto de evaluación de información disponible durante el entrenamiento, lo que sobreestima el desempeño.}
}

\newglossaryentry{corpuslimpio}
{
    name={Corpus limpio},
    description={Conjunto de payloads de evaluación filtrado para excluir los que coinciden con el entrenamiento; en este trabajo es el de la Ronda~5.}
}

\newglossaryentry{useragent}
{
    name={\texttt{User-Agent}},
    sort={User-Agent},
    description={Cabecera HTTP con la que el cliente se identifica (navegador, cURL, entre otros); su representación en el corpus sesgó al modelo.}
}

\newglossaryentry{falsopositivo}
{
    name={Falso positivo},
    description={Petición legítima que el sistema clasifica o bloquea como ataque.}
}

\newglossaryentry{parmodelos}
{
    name={Par de modelos},
    description={Combinación de un archivo \texttt{rf.onnx} y uno \texttt{if.onnx} identificada por su suma SHA-256 (pares A, B y C).}
}

\newglossaryentry{tarball}
{
    name={\textit{tarball}},
    sort={tarball},
    description={Archivo comprimido \texttt{.tgz} de un paquete de npm, publicado o construido localmente.}
}

\newglossaryentry{kanban}
{
    name={Kanban},
    description={Metodología de gestión basada en un tablero de columnas que representan el estado de las tareas.}
}

\newglossaryentry{artillery}
{
    name={Artillery},
    description={Herramienta de pruebas de carga usada para medir la latencia con y sin el middleware.}
}

\newglossaryentry{sqlmap}
{
    name={sqlmap},
    description={Herramienta de código abierto que automatiza la detección y la explotación de inyecciones SQL; generó los payloads de SQL Injection del corpus limpio.}
}

\newglossaryentry{zap}
{
    name={OWASP ZAP},
    description={Herramienta de OWASP para pruebas de seguridad de aplicaciones web; se usó para confirmar el payload de XSS.}
}

\newglossaryentry{mcp}
{
    name={MCP},
    description={Protocolo de Contexto de Modelo (\textit{Model Context Protocol}): estándar para que un asistente de inteligencia artificial use herramientas externas.}
}

\fi

\usepackage[top=1in, left=1.5in, right=1in, bottom=1in]{geometry}
\usepackage[Sonny]{fncychap}

\graphicspath{ {figuras/} }
\definecolor{uvg-green}{RGB}{17,71,52}
\newcommand{\defaultparformat}[1]{
	{\setlength{\parskip}{2ex}
     \input{#1}}
}
\ifdefined\capsecuvg
	\renewcommand\thechapter{\Roman{chapter}}
    \renewcommand\thesection{\Alph{section}}
	\renewcommand\thesubsection{\arabic{subsection}}
    \renewcommand\thesubsubsection{\alph{subsubection}}
\fi
\counterwithout{figure}{chapter}
\counterwithout{table}{chapter}
\counterwithout{equation}{chapter}

\providecommand{\blankpage}{
\newpage
\thispagestyle{empty}
\mbox{}
\newpage
}

% ==============================================================================
% Paquetes y comandos adicionales del usuario
% ==============================================================================
% TikZ: requerido por las figuras del capítulo de Marco Teórico (pipeline de
% middleware, comparación RASP vs. control perimetral).
\usepackage{tikz}
\usetikzlibrary{arrows.meta, positioning, shapes.geometric, fit, backgrounds}


\pagestyle{headings}
\begin{document}

\ifdefined\parpordefecto
	\setlength{\parskip}{2ex}
\fi

\fi

\chapter{Metodología}
\label{cap:metodologia}

Este capítulo describe cómo se ejecutó el componente de ciberseguridad de \textit{logSguarDian} y
cómo se obtuvieron, procesaron y validaron los datos que sustentan el capítulo de Resultados.
Constituye una evolución del Marco Metodológico del Protocolo de Graduación: conserva el
enfoque, la metodología de gestión, las seis fases y las métricas allí definidas, e incorpora los
cambios que el avance real del proyecto hizo necesarios. Cada cambio se declara de forma
explícita en la Sección~\ref{sec:met-cambios}, junto con su causa; ninguna cifra o criterio del
protocolo se sustituye en silencio.

El proyecto integrador se compone de dos componentes desarrollados en paralelo. Este documento
cubre el componente de ciberseguridad; el diseño, entrenamiento y serialización de los modelos
de aprendizaje automático pertenecen al componente de Inteligencia Artificial y aquí se tratan
únicamente como entradas, con su contrato de interfaz (Sección~\ref{subsec:met-f2}).

% ------------------------------------------------------------------------------
\section{Enfoque y tipo de investigación}
\label{sec:met-enfoque}

El trabajo adopta un enfoque cuantitativo, de tipo aplicado y con diseño experimental. Es
cuantitativo porque todas las preguntas que plantean los objetivos se responden mediante
mediciones numéricas sobre conjuntos de datos y condiciones de carga definidos de antemano,
contrastadas con criterios de aceptación fijados antes de la medición
\cite{hernandez2014}. Es aplicado porque su producto, una librería npm publicada, responde a un
problema operativo concreto: la ausencia de protección a nivel de aplicación en servicios
Node.js de organizaciones sin equipo de seguridad dedicado.

El diseño es experimental en el sentido de que el investigador manipula de forma deliberada la
variable independiente y mide su efecto sobre variables dependientes en un entorno controlado
\cite{hernandez2014}. La variable independiente es la configuración de defensa del servidor
anfitrión, con cuatro niveles: sin protección (Configuración~1, grupo de control), con
\textit{logSguarDian} (Configuración~2), con un WAF ModSecurity con OWASP CRS (Configuración~3a)
y con ambas capas superpuestas (Configuración~3b). Las variables dependientes son la tasa de
detección por categoría de ataque, la tasa de falsos positivos sobre tráfico legítimo, la
latencia p95 de respuesta y el consumo de memoria y CPU. La línea base sin middleware y el
servidor anfitrión idéntico entre condiciones (mismo árbol de la aplicación, mismo
\texttt{Dockerfile}, misma base de datos sembrada) controlan las variables extrañas que podrían explicar las diferencias observadas.

La unidad de análisis es la petición HTTP individual dirigida a un endpoint de una aplicación
Express, caracterizada por su método, ruta, cadena de consulta, cuerpo y cabeceras. No participan
personas como sujetos de estudio; por esa razón no aplica un muestreo de población humana, y la
Sección~\ref{sec:met-unidad} define las poblaciones y muestras de peticiones utilizadas.

% ------------------------------------------------------------------------------
\section{Metodología de trabajo}
\label{sec:met-trabajo}

\subsection{Kanban como metodología de gestión}
La gestión del trabajo se realizó con Kanban \cite{anderson2010}, sobre un tablero en Notion
compartido por los dos integrantes del proyecto. La selección, justificada en el protocolo, se
mantuvo por tres razones que el desarrollo confirmó en la práctica:

\begin{itemize}
    \item \textbf{Ausencia de un receptor externo de incrementos.} No existe un Product Owner
    ni un cliente que valide entregas periódicas; las ceremonias de Scrum habrían carecido de
    propósito metodológico real.
    \item \textbf{Dependencias técnicas secuenciales.} El corpus es prerrequisito del
    entrenamiento, y el modelo serializado es prerrequisito de la integración en el middleware.
    Kanban representa ese flujo con columnas de estado sin imponer iteraciones de duración fija
    que obliguen a avanzar con una dependencia sin resolver.
    \item \textbf{Respuesta continua a hallazgos experimentales.} Durante el proyecto aparecieron
    hallazgos que obligaron a reabrir trabajo ya cerrado: la fuga de datos entre el corpus de
    evaluación y el de entrenamiento, el sesgo del modelo hacia la cabecera \texttt{User-Agent} y
    el colapso del servidor bajo volumen de datos realista. Kanban permitió repriorizarlos de
    inmediato sin esperar al cierre de un sprint.
\end{itemize}

El tablero se organizó en cuatro columnas que representan los estados posibles de una tarea:
No empezado, En proceso, En revisión y Completado. Respecto al protocolo, el tablero no incluyó
una columna Bloqueado. El criterio de salida de la columna En revisión es la verificación del indicador
cuantitativo asociado a la tarea (Sección~\ref{sec:met-metricas}), no una apreciación subjetiva
de completitud.

\subsection{Flujo de control de versiones e integración continua}
\label{subsec:met-flujo}
El tablero se complementó con un flujo de ramas en Git que convierte la revisión en un control
obligatorio y no opcional:

\begin{itemize}
    \item \textbf{Ramas.} Todo cambio se desarrolla en una rama propia con prefijo por tipo
    (\texttt{feat/}, \texttt{fix/}, \texttt{test/}, \texttt{docs/}, \texttt{perf/},
    \texttt{chore/}) creada desde \texttt{develop}. La integración sigue la secuencia
    \texttt{rama} $\rightarrow$ \texttt{develop} $\rightarrow$ \texttt{main}; un flujo de
    GitHub Actions rechaza cualquier pull request hacia \texttt{main} cuyo origen no sea
    \texttt{develop}.
    \item \textbf{Aprobación restringida.} Las reglas de protección de rama del repositorio solo
    permiten que los dos integrantes del proyecto aprueben y fusionen pull requests, de modo que
    cada cambio de un integrante es revisado por el otro.
    \item \textbf{Integración continua (CI).} Cada pull request ejecuta en GitHub Actions
    (\texttt{ubuntu-latest}) los trabajos de análisis estático (ESLint), compilación
    (\texttt{tsc}), pruebas del extractor de características (incluida la prueba de paridad
    numérica contra la implementación de referencia en Python), pruebas del paquete
    \texttt{core} (middleware, almacenamiento, webhooks y la interfaz de línea de comandos, CLI) y pruebas del paquete
    \texttt{mlops}. La fusión de \texttt{develop} a \texttt{main} publica automáticamente el
    paquete en el registro npm si la versión no ha sido publicada antes.
    \item \textbf{Trazabilidad de decisiones.} Las decisiones técnicas y su justificación se
    registran en una bitácora versionada, con una entrada por tarea cerrada, y cada
    reentrenamiento del modelo va precedido de un commit que fija el hash del conjunto de prueba
    (Sección~\ref{subsec:met-f1}), con una excepción: el rebloqueo del 23 de septiembre se hizo
    después de restaurar los datos de la generación vigente.
\end{itemize}

Este flujo tiene una finalidad metodológica, no solo de organización: cada cifra reportada en
Resultados queda asociada a un commit, a un script versionado y, cuando aplica, a una ejecución
de GitHub Actions cuyo registro puede consultarse posteriormente.

\subsection{Distribución de roles}
El Cuadro~\ref{tab:met-roles} distribuye las responsabilidades. Las responsabilidades del
componente de IA se listan solo para delimitar la frontera entre componentes.

\begin{table}[H]
\centering
\small
\begin{tabular}{|L{2.6cm}|L{4.9cm}|L{5.3cm}|}
\hline
\textbf{Integrante} & \textbf{Rol principal} & \textbf{Responsabilidades clave} \\ \hline
Sebastián Huertas & Componente de ciberseguridad & Extractor inicial de 72 características y el
paso a \texttt{worker\_threads}; selección y mapeo de vectores de ataque;
auditoría de fuentes y etiquetado; entorno Docker de evaluación; capa de captura y respuesta
(HTTP 403, registro, webhooks, CLI); calibración del umbral de bloqueo; corpus de evaluación
limpio; pruebas de penetración controlada; benchmarks de latencia, memoria y CPU; validación en
aplicaciones de terceros. \\ \hline
Diego Valenzuela & Componente de IA & Análisis y justificación de las características, definición de
las no disponibles en inferencia y del contrato de entrada; cuatro características adicionales y
correcciones del extractor; diseño de la arquitectura de inferencia en hilos (hilo dedicado de Random
Forest, grupo de Isolation Forest y lectura asíncrona de su resultado); entrenamiento, selección y
serialización ONNX de los modelos; pipeline de reentrenamiento offline (\texttt{mlops}). \\ \hline
Ambos & Integración y documentación & Arquitectura de la librería; pruebas de integración
extremo a extremo; revisión cruzada de pull requests; publicación en npm; documentación técnica
final. \\ \hline
\end{tabular}
\caption[Distribución de roles y responsabilidades]{Distribución de roles y responsabilidades del equipo}
\label{tab:met-roles}
\end{table}
{\footnotesize Nota. Elaboración propia.}

La frontera entre componentes en el pipeline de características es la siguiente. El extractor
inicial de 72 características y el paso a \texttt{worker\_threads} los
implementó Sebastián Huertas. Diego Valenzuela diseñó la arquitectura de inferencia en hilos (hilo
dedicado de Random Forest, grupo de Isolation Forest y lectura asíncrona de su resultado), analizó
y justificó cada característica, definió cuáles no existen en inferencia y fijó el contrato de
entrada, añadió cuatro características y las correcciones del extractor surgidas del análisis del
modelo, y desarrolló el entrenamiento, la serialización ONNX y el pipeline de reentrenamiento.

% ------------------------------------------------------------------------------
\section{Unidad de estudio, población y muestras}
\label{sec:met-unidad}

La población de interés es el conjunto de peticiones HTTP que puede recibir un endpoint de una
aplicación Node.js/Express, legítimas o pertenecientes a alguna de las cuatro categorías de
ataque del Objetivo Específico~1. Como esa población no es enumerable, el estudio emplea cinco
muestras, cada una construida para responder una pregunta distinta. El Cuadro~\ref{tab:met-muestras}
las resume.

\begin{table}[H]
\centering
\small
\begin{tabular}{|L{2.8cm}|L{3.6cm}|L{2.4cm}|L{3.9cm}|}
\hline
\textbf{Muestra} & \textbf{Composición} & \textbf{Tipo de muestreo} & \textbf{Pregunta que responde} \\ \hline
M1. Corpus etiquetado & 383{,}202 peticiones de siete fuentes públicas auditadas, dos listas públicas de payloads, conjuntos sintéticos y telemetría curada:
268{,}240 de entrenamiento, 57{,}481 de validación y 57{,}481 de prueba & Estratificado por
clase y fuente (70/15/15) & ¿El modelo clasifica correctamente las cuatro categorías? (OE2, OE3) \\ \hline
M2. Corpus limpio de evaluación (Ronda~5) & 200 SQLi (sqlmap), 200 Path Traversal y 200 Command
Injection (generadores propios), y 200 XSS (generador propio, 25 de octubre de 2026); el payload único de XSS confirmado con OWASP ZAP se conserva aparte como hallazgo & Intencional, con filtro de exclusión contra el entrenamiento &
¿Qué proporción de payloads nunca vistos detecta el sistema desplegado? (OE3) \\ \hline
M3. Tráfico de carga & Navegación legítima y tráfico mixto generados con Artillery a
$\approx$100 peticiones/s & Generado, cinco repeticiones por condición & ¿Cuánto degrada el
middleware el rendimiento del servidor? (Objetivo general) \\ \hline
M4. Tráfico de navegador & Rutas sin parámetros (\texttt{/login}, \texttt{/posts},
\texttt{/profile}) con cadenas \texttt{User-Agent} de navegadores reales & Intencional & ¿El
modelo bloquea tráfico legítimo de navegadores reales? (OE3) \\ \hline
M5. Aplicaciones de terceros & DVNA \cite{appsecco2026dvna}, node-api-goat \cite{layro012026nodeapigoat}, dvws-node \cite{snoopysecurity2026dvwsnode} y OWASP Juice Shop \cite{kimminich2026juiceshop} & Intencional,
por diversidad de arquitectura & ¿La calibración generaliza fuera de la aplicación de
referencia? (Objetivo general) \\ \hline
\end{tabular}
\caption[Muestras de peticiones utilizadas en el estudio]{Muestras de peticiones HTTP utilizadas en el estudio}
\label{tab:met-muestras}
\end{table}
{\footnotesize Nota. Elaboración propia con datos de los repositorios del proyecto.}

La muestra M1 es un insumo compartido con el componente de IA; su tamaño satisface el mínimo de
100{,}000 muestras del OE2. La muestra M2 reemplaza al corpus de evaluación inicial (Ronda~4),
descartado por fuga de datos (Sección~\ref{subsec:met-fuga}). La categoría XSS de M2 se amplió el 25 de octubre de 2026 con un corpus propio de 200 payloads
(Procedimiento~4.3); el payload único confirmado con OWASP ZAP se conserva como hallazgo de
explotabilidad (n=1), con la restricción declarada en la Sección~\ref{sec:met-restricciones}.

% ------------------------------------------------------------------------------
\section{Fases, actividades y procedimientos}
\label{sec:met-fases}

El proyecto conserva las seis fases del protocolo. El Cuadro~\ref{tab:met-fases} compara el
período planificado con el período real, reconstruido a partir del historial de commits de los
repositorios del proyecto. Las fases se solaparon más de lo previsto: la detección de problemas
durante las pruebas obligó a reabrir las fases de construcción del dataset y de entrenamiento,
lo que corresponde al ciclo iterativo descrito en el protocolo.

\begin{table}[H]
\centering
\small
\begin{tabular}{|L{2.6cm}|L{1.9cm}|L{2.1cm}|L{1.2cm}|L{4.6cm}|}
\hline
\textbf{Fase} & \textbf{Protocolo} & \textbf{Real} & \textbf{OE} & \textbf{Entregable verificable} \\ \hline
F1. Construcción del dataset & 1--31 may & 11 may -- 20 jun & OE1, OE2 & Corpus $\geq$100{,}000
muestras etiquetadas; auditoría de fuentes y licencias; taxonomía de cinco clases \\ \hline
F2. Entrenamiento del modelo (IA) & 1--19 jun & 11 jun -- 26 sep & OE3 & Modelos ONNX con paridad
$<$0.1~\%; lectura del conjunto de prueba \\ \hline
F3. Desarrollo de la librería & 20 jun -- 1 ago & 30 jun -- 30 ago & OE3 & Paquete
\texttt{logsguardian@0.1.0} publicado en npm, con cuatro grupos de comandos CLI probados con
Jest \\ \hline
F4. Pruebas & 2 ago -- 30 sep & 23 jul -- 10 oct & OE3 & Resultados de las configuraciones 1, 2,
3a y 3b; corpus limpio; benchmark de rendimiento; validación en aplicaciones de terceros \\ \hline
F5. Aplicación de feedback & 15 ago -- 30 sep & 24 jul -- 11 oct & OE3 & Registro de iteraciones
con métricas antes y después de cada corrección; versión 0.1.1 publicada en npm \\ \hline
F6. Documentación & octubre & 16 ago -- oct & Todos & Informe final; README publicado en npm;
documentación técnica \\ \hline
\end{tabular}
\caption[Fases del proyecto: período planificado y real]{Fases del proyecto: período planificado en el protocolo, período real de ejecución y entregables verificables}
\label{tab:met-fases}
\end{table}
{\footnotesize Nota. Elaboración propia. El período real corresponde al primer y último commit
relacionado con cada fase en los repositorios \texttt{logSguarDian} y
\texttt{logSguarDian-vulnerable-project}; las fechas de publicación en npm son las del registro
(UTC). Todas las fechas son de 2026.}

Las subsecciones siguientes describen los procedimientos de cada fase. Para cada uno se indica
su finalidad y el objetivo al que contribuye.

\subsection{Fase 1: análisis de vectores y construcción del dataset}
\label{subsec:met-f1}

\textbf{Procedimiento 1.1. Mapeo de vectores a marcos de referencia (OE1).} Cada uno de los
cuatro vectores se asoció a su categoría del OWASP Top 10 2021, a su debilidad CWE y a su
técnica MITRE ATT\&CK, verificando cada correspondencia contra la página oficial del marco
\cite{owasp2021, owaspa01_2021, owaspa03_2021, mitret1190, mitret1189}. El mapeo distingue
SQL Injection, XSS y Command Injection (A03:2021, Injection) de Path Traversal/LFI (A01:2021,
Broken Access Control). Finalidad: fijar las cinco clases canónicas del dataset (\texttt{sqli},
\texttt{xss}, \texttt{path\_traversal}, \texttt{cmdi}, \texttt{benign}) y las categorías sobre
las que se calculan las métricas del OE3.

\textbf{Procedimiento 1.2. Auditoría de fuentes (OE2).} Se auditaron nueve fuentes públicas
candidatas, registrando para cada una su licencia (verificada contra la API o la página del
repositorio de origen), su DOI cuando existe, su formato, su número de registros y las
categorías que cubre. Se aprobaron ocho y se excluyó el dataset AIT-LDSv2 \cite{landauer2022ait,landauer2023maintainable} porque sus etiquetas
describen el rol del atacante y no el tipo de payload, lo que habría exigido un reetiquetado manual
con riesgo de error superior a su aporte de 3{,}435 filas. De las ocho aprobadas, siete aportan
filas al corpus; la lista de payloads de Path Traversal de omurugur \cite{ugur2019traversal} se aprobó, pero no tiene
analizador y no aporta filas. El corpus final incluye además dos listas públicas de payloads de Command Injection (SecLists y PayloadsAllTheThings
\cite{miessler2026seclists,swisskyrepo2026payloads}), incorporadas después de la auditoría y cuya
licencia (MIT, consultada en sus repositorios el 10 de octubre de 2026) no consta en ella, siete
conjuntos sintéticos propios y 119 filas de telemetría curada del pipeline de reentrenamiento
(Cuadro~\ref{tab:met-origenes}). El dataset de peticiones bloqueadas por OWASP ModSecurity
publicado en Zenodo \cite{lucz2025,lucz2025data}, que el protocolo designaba como base, aporta el 1.0~\% de las
filas; el volumen proviene principalmente de la colección CAPEC
\cite{sureda2022srbh,sureda2022newmultilabel} (55.7~\%) y de modsec-learn
\cite{pralab2024modseclearn,scano2024modseclearn} (26.0~\%). Las fuentes de licencia no especificada por su origen (XSS\_dataset.csv y command
injection.csv, dos de las siete auditadas presentes en el corpus, 3.4~\% de las filas) no se
redistribuyen: sus archivos de trabajo no se versionan en el repositorio y el proyecto entrega
únicamente los scripts de procesamiento y los vectores de características derivados.

\begin{table}[H]
\centering
\small
\begin{tabular}{|L{4.3cm}|L{1.8cm}|L{1.3cm}|L{5.3cm}|}
\hline
\textbf{Origen} & \textbf{Filas} & \textbf{\%} & \textbf{Licencia} \\ \hline
Colección CAPEC (SR-BH 2020) \cite{sureda2022srbh,sureda2022newmultilabel} & 213{,}514 & 55.72 & CC0-1.0, auditada \\ \hline
modsec-learn \cite{pralab2024modseclearn,scano2024modseclearn} & 99{,}645 & 26.00 & MIT, auditada \\ \hline
payload\_full.csv \cite{thakur2025payloads} & 29{,}719 & 7.76 & MIT, auditada \\ \hline
Payloads.csv \cite{shahane2021xss} & 18{,}411 & 4.80 & CC BY 4.0, auditada \\ \hline
XSS\_dataset.csv \cite{shah2020xss} & 10{,}854 & 2.83 & No especificada por el origen, auditada \\ \hline
OWASP ModSecurity (Zenodo) \cite{lucz2025} & 3{,}860 & 1.01 & CC BY 4.0, auditada \\ \hline
command injection.csv \cite{pawase2024cmdi} & 2{,}053 & 0.54 & No especificada por el origen, auditada \\ \hline
SecLists y PayloadsAllTheThings \cite{miessler2026seclists,swisskyrepo2026payloads} (Command Injection) & 2{,}470 & 0.64 & MIT; no consta en la auditoría \\ \hline
Conjuntos sintéticos propios (siete) & 2{,}557 & 0.67 & Propios \\ \hline
Telemetría curada del pipeline de reentrenamiento & 119 & 0.03 & Propia \\ \hline
\textbf{Total} & \textbf{383{,}202} & \textbf{100.00} & \\ \hline
\end{tabular}
\caption[Orígenes de las filas del corpus]{Orígenes de las filas del corpus etiquetado y estado de su licencia}
\label{tab:met-origenes}
\end{table}
{\footnotesize Nota. Elaboración propia con el conteo de la columna \texttt{\_source} del archivo de
vectores de características (repositorio \texttt{logSguarDian}, 10 de octubre de 2026) y con
\texttt{docs/dataset-audit.md}.}

\textbf{Procedimiento 1.3. Normalización y etiquetado determinista (OE2).} Un analizador por
fuente transforma cada registro crudo a un esquema canónico común (\texttt{CanonicalRequest}:
método, ruta, consulta, cuerpo y cabeceras), y un archivo de mapeo versionado
(\texttt{label\_map.yaml}) asigna la clase canónica a partir de la etiqueta original de cada
fuente. El etiquetado no involucra juicio manual por muestra: la misma entrada produce siempre la
misma etiqueta. Las decisiones de mapeo que no son equivalencias conceptuales se declaran como
tales; por ejemplo, la asignación de CAPEC-242 (Code Injection) a la clase \texttt{xss}.

\textbf{Procedimiento 1.4. Deduplicación.} Se eliminan los registros idénticos mediante un hash
SHA-256 de la concatenación \texttt{ruta|consulta|cuerpo|etiqueta}. Los casi-duplicados
se detectan de dos formas y se reportan, pero no se eliminan: por distancia de Levenshtein \cite{levenshtein1966binary} menor
que 3 sobre la consulta, dentro de la misma clase, y por MinHash \cite{broder1997resemblance} con LSH \cite{indyk1998approximate} (128 permutaciones y
bandas de 8 por 16), que se aplicó a las cinco clases. La decisión de marcar sin eliminar se
declara como amenaza a la validez (Sección~\ref{sec:met-validez}).

\textbf{Procedimiento 1.5. Partición y bloqueo del conjunto de prueba (OE2, OE3).} El corpus
unificado se divide 70/15/15 en entrenamiento, validación y prueba, con estratificación conjunta
por clase y fuente y semilla fija. Antes de escribir las particiones, el script verifica tres
condiciones: cero solapamiento de hashes entre particiones, distribución de clases por partición
dentro de $\pm$2~\% de la global y distribución de fuentes dentro de $\pm$5~\%. Si se cumplen,
calcula un hash SHA-256 sobre los hashes ordenados de las filas de prueba y lo versiona
(\texttt{test.lock.sha256}) antes de cualquier entrenamiento. Finalidad: que cualquier lectura
posterior del conjunto de prueba pueda demostrar que la partición leída es la que se bloqueó.

\subsection{Fase 2: entrenamiento del modelo (entrada del componente de IA)}
\label{subsec:met-f2}

El entrenamiento de un Random Forest como clasificador supervisado y de un Isolation Forest como
detector de anomalías, ambos con scikit-learn \cite{pedregosa2011}, corresponde al componente de
IA. Para este componente se fijó el siguiente contrato de interfaz, que es lo que el middleware
consume:

\begin{itemize}
    \item Vector de características calculado por una única implementación en TypeScript
    (\texttt{packages/extractor}), consumida tanto por el pipeline de entrenamiento como por el
    middleware, para eliminar la divergencia entre entrenamiento y producción. El vector tiene 76
    características. En la generación vigente (\texttt{rf\_v11}/\texttt{if\_v10}), el Random Forest
    consume 69 y el Isolation Forest 63; en la generación publicada (\texttt{rf\_v10}/\texttt{if\_v9}),
    67 y 61. Se excluyen por nombre las que no están disponibles en el momento de inspeccionar la
    petición (código de estado y cinco características temporales) y
    \texttt{non\_json\_quote\_count}, agregada al extractor después del entrenamiento de los
    modelos vigentes. El Isolation Forest excluye además seis características con varianza nula o
    casi nula en tráfico legítimo.
    \item Modelos serializados en ONNX y ejecutados con \texttt{onnxruntime-node}
    \cite{onnxruntimedocs}, con una prueba de paridad obligatoria (diferencia máxima entre las
    predicciones de Python y de Node.js menor que 0.1~\%) antes de incorporarlos al paquete.
    \item Cada evaluación se ata a un par concreto de binarios ONNX, identificado por su hash
    SHA-256. Durante el proyecto se usaron tres pares, descritos en el
    Cuadro~\ref{tab:pares-modelos}.
\end{itemize}

\begin{table}[H]
\centering
\small
\begin{tabular}{|L{0.9cm}|L{2.6cm}|L{1.4cm}|L{2.6cm}|L{5.2cm}|}
\hline
\textbf{Par} & \textbf{Versión} & \textbf{Caract. (RF/IF)} & \textbf{Hash (\texttt{rf.onnx} / \texttt{if.onnx})} & \textbf{Evaluaciones en que se usó} \\ \hline
A & \texttt{rf\_v10} / \texttt{if\_v9} (2 de agosto de 2026) & 67 / 61 & \texttt{18083725...} / \texttt{5de374a7...} & Paquete \texttt{logsguardian@0.1.0} (commit \texttt{c3d591a}); Ronda~4 (Configuraciones 2, 3a y 3b sobre los 590 payloads)\textsuperscript{c}; medición principal de la Ronda~5; estudio de \texttt{User-Agent}\textsuperscript{a}; DVNA, node-api-goat y dvws-node (atribución probable)\textsuperscript{b} \\ \hline
B & \texttt{rf\_v11} / \texttt{if\_v10}, primer build (30 de agosto de 2026) & 69 / 63 & \texttt{a7c015f5...} / \texttt{3cf8e45c...} & Repetición de la Configuración~2 sobre los 590 payloads de la Ronda~4 (9 de septiembre de 2026); benchmark de latencia, memoria y CPU en Linux nativo \\ \hline
C & \texttt{rf\_v11} / \texttt{if\_v10}, reentrenamiento (26 de septiembre de 2026) & 69 / 63 & \texttt{25b407e6...} / \texttt{5be28ee8...} & Paquete \texttt{logsguardian@0.1.1} (publicado el 11 de octubre de 2026)\textsuperscript{d}, con el que se verificó la cobertura de Path Traversal y de Command Injection de la Ronda~5 con el ejecutor corregido; lectura del conjunto de prueba, paridad ONNX y análisis SHAP; repetición de la Ronda~5 con el tarball de validación 22550ea; OWASP Juice Shop (copia local); Configuraciones 2 y 3b de la medición unificada del 25 de octubre de 2026 (paquete 0.1.1 del registro, archivo \texttt{abf8d5c4...}) \\ \hline
\end{tabular}
\caption[Pares de modelos ONNX y su uso]{Pares de modelos ONNX utilizados y evaluaciones asociadas a cada uno}
\label{tab:pares-modelos}
\end{table}
{\footnotesize Nota. Elaboración propia con datos de \texttt{docs/model-provenance.md} (repositorio
\texttt{logSguarDian}), verificados por hash contra el historial de git y contra el paquete
descargado del registro de npm. \textsuperscript{a}El archivo de resultados no registra el modelo; la
atribución se apoya en el contenido de \texttt{models/} en el commit del estudio (\texttt{5966712}),
que contiene el par~A (hash verificado). \textsuperscript{b}Los tres proyectos declaran la dependencia
\texttt{\^{}0.1.0} del registro de npm, cuyo único contenido publicado en la fecha de las mediciones
(30 de septiembre de 2026) era el par~A, pero el paquete y el modelo exactos de cada medición no
constan en el repositorio. \textsuperscript{c}En las Rondas 1 a 3,
con un corpus pequeño, el modelo no consta en los resultados. \textsuperscript{d}Contenido verificado
descargando el paquete del registro: modelos con los hashes del par~C, extractor de 76
características, corrección de la recodificación de la consulta y umbral del Isolation Forest leído del archivo \texttt{parity\_report.json} del modelo (0.004206). No se agrega un par~D: el paquete 0.1.1 y el tarball
de validación contienen los mismos modelos, y los pares se identifican por el hash de sus archivos
ONNX; difieren en el middleware, no en el modelo.}

El paquete \texttt{logsguardian@0.1.0} se publicó en npm el 30 de agosto de 2026 desde el commit
\texttt{c3d591a} y contiene el par~A; el primer build de \texttt{rf\_v11}/\texttt{if\_v10} (par~B) se
generó ese mismo día y no se incluyó en él. La versión \texttt{0.1.1}, publicada el 11 de octubre de
2026 (SHA-256 del archivo \texttt{abf8d5c4...}), contiene el par~C, el extractor de 76
características, la corrección del defecto de recodificación de la consulta y el umbral del Isolation Forest leído del archivo \texttt{parity\_report.json} del modelo. En este
documento, \textit{paquete 0.1.0}
designa a la primera y \textit{paquete 0.1.1} a la segunda. Con la 0.1.1 se midió la cobertura de la Configuración~2 de la Ronda~5, con el ejecutor
original y con el corregido, y, el 25 de octubre de 2026, la de las Configuraciones 2 y 3b con el
ejecutor corregido; las aplicaciones de terceros y el rendimiento no se midieron con ella,
que es posterior a esas pruebas. Las mediciones de la Ronda~5 designadas como del paquete 0.1.0
se hicieron con un tarball construido en el commit \texttt{a3aada0} (SHA-256 \texttt{e09120006dd1...}),
que contiene el par~A pero no es el archivo publicado en el registro (SHA-256
\texttt{c5c1a1da57cf...}); se les llama paquete 0.1.0 por contener el mismo par de modelos.
Para medir la cobertura del modelo vigente antes de que existiera la 0.1.1, se construyó un
paquete a partir del commit \texttt{22550ea} de la rama \texttt{develop}, con el par~C y sin la
corrección, denominado \textit{tarball de validación 22550ea} (SHA-256 \texttt{cf1238de...}). Ese
paquete no se publicó y no debe confundirse con la 0.1.0 ni con la 0.1.1. La evaluación de
cobertura sobre el corpus limpio se ejecutó con el paquete 0.1.0 y con el tarball de validación,
de modo que esa métrica existe para los pares A y C, y para la versión 0.1.1 publicada en las Configuraciones 2 y 3b.

El pipeline de reentrenamiento offline con compuertas automáticas (\texttt{packages/mlops}) se
incorporó por recomendación del asesor como parte del componente de IA. No modifica el
comportamiento de la librería publicada: los modelos se entrenan fuera de línea y se distribuyen
fijos con cada versión, sin aprendizaje a partir del tráfico del usuario.

\textbf{Procedimiento 2.1. Calibración del umbral de bloqueo (OE3, ciberseguridad).} La
política de decisión otorga al Random Forest la autoridad exclusiva de bloqueo: una petición se
bloquea cuando la probabilidad de la clase de ataque más probable es igual o mayor que el umbral
\texttt{RF\_THRESHOLD}. El Isolation Forest nunca bloquea; solo marca la petición como anómala
en el registro (\texttt{pass\_anomaly}), porque una política disyuntiva (bloquear si cualquiera
de los dos modelos alerta) habría bloqueado cerca del 10~\% del tráfico legítimo. El umbral se
calibró con el modelo \texttt{rf\_v3} mediante un barrido entre 0.10 y 0.70 sobre la partición
de validación, nunca sobre la de prueba, registrando para cada valor la precisión, el recall, los
falsos positivos y la detección por clase. Se seleccionó 0.35, el extremo conservador de la
meseta en que recall y falsos positivos permanecen constantes. El repositorio registra una
confirmación de ese umbral sobre el conjunto de prueba de \texttt{rf\_v3}, pero no conserva el
artefacto de esa lectura, por lo que esa confirmación no es verificable. Para \texttt{rf\_v11} el
umbral se mantuvo sin recalibrar. El umbral del Isolation Forest se seleccionó también sobre
validación y se fija en los metadatos de cada modelo; sin embargo, en dos generaciones previas
(\texttt{if\_v2} e \texttt{if\_v5}) el objetivo interno de calibración (falsos positivos de 8~\% y
luego de 6~\%) se definió después de observar un resultado de prueba que no lo cumplía.

\textbf{Procedimiento 2.2. Lectura del conjunto de prueba (OE3).} La evaluación final de los
modelos publicados ejecuta un script versionado (\texttt{evaluate\_test.py}) que: (1) recalcula el
hash del conjunto de prueba y aborta si no coincide con el bloqueado; (2) carga los binarios ONNX
del repositorio (par C), no los modelos de Python; (3) calcula el reporte de clasificación y el AUC-ROC
por clase en esquema uno contra el resto, asignando a cada petición la clase de mayor
probabilidad (sin aplicar \texttt{RF\_THRESHOLD}); (4) evalúa el Isolation Forest con el umbral de
sus metadatos; y (5) guarda un manifiesto de sumas de verificación del conjunto de prueba y de
ambos binarios junto al resultado. Finalidad: que cada cifra del OE3 quede atada al contenido
exacto, byte a byte, del modelo y de los datos que la produjeron.

Para la generación \texttt{rf\_v11} el conjunto de prueba se leyó tres veces: el 24 de septiembre
de 2026, con un modelo entrenado antes del rebloqueo del 23 de septiembre y un umbral del
Isolation Forest de 0.00807; y el 25 y el 26 de septiembre, con modelos reentrenados sobre la
partición vigente y un umbral de 0.00421. Las cifras reportadas en Resultados corresponden a la
lectura del 26 de septiembre (par C).

\subsection{Fase 3: desarrollo de la librería}
\label{subsec:met-f3}

\textbf{Procedimiento 3.1. Captura y respuesta.} El middleware transforma el objeto
\texttt{req} de Express al esquema canónico, extrae el vector de características y envía la
inferencia a \texttt{worker\_threads}, para no bloquear el Event Loop. La arquitectura de esa
inferencia en hilos (un hilo dedicado al Random Forest, un grupo rotativo de dos hilos para el
Isolation Forest, y una respuesta que se resuelve con el veredicto del Random Forest sin esperar
al Isolation Forest, cuya respuesta tardía actualiza el registro de forma asíncrona) es una
decisión de diseño del componente de Inteligencia Artificial, justificada en el trabajo de
graduación de Diego Valenzuela \cite{valenzuela2026ia}; aquí se describe solo como la interfaz que el middleware consume.
Ante una petición bloqueada, el middleware responde HTTP 403 en
modo \texttt{block}; en modo \texttt{monitor} solo registra. Todo evento se persiste en SQLite y
los eventos accionables (\texttt{block}, \texttt{pass\_anomaly}) se notifican a los webhooks
registrados. Ante un error o un tiempo de espera excedido del hilo de inferencia (50~ms por
defecto), el middleware deja pasar la petición (\textit{fail-open}), para que un fallo interno
no convierta la librería en un vector de denegación de servicio contra la aplicación anfitriona.

\textbf{Procedimiento 3.2. Verificación funcional.} Cada módulo se valida con pruebas
automatizadas en Jest \cite{jestdocs} y Supertest \cite{supertestdocs} que se ejecutan en
integración continua. Las correcciones de defectos se acompañan de una prueba de regresión y se
verifican mediante sabotaje y reversión: se reintroduce el defecto de forma temporal, se
confirma que la prueba nueva falla y se revierte. Antes de la publicación, el paquete empaquetado
(\texttt{pnpm pack}) se instaló en un directorio limpio fuera del monorepo y se probó contra una
aplicación Express real con una petición SQLi (respuesta 403) y una legítima (respuesta 200).

\subsection{Fase 4: pruebas}
\label{subsec:met-f4}

Las pruebas se ejecutan sobre una aplicación de referencia deliberadamente vulnerable
(\texttt{logSguarDian-vulnerable-project}), construida con Express, PostgreSQL y plantillas EJS,
y desplegada con Docker Compose \cite{dockerdocs, moreau2023}. La aplicación expone endpoints
vulnerables a cada vector: \texttt{GET /posts/search?q=} y \texttt{POST /login} (SQLi),
\texttt{POST /posts} y \texttt{POST /profile} (XSS), \texttt{GET /posts/:id/attachment?file=}
(Path Traversal/LFI) y \texttt{POST /admin/ping} (Command Injection).

\textbf{Procedimiento 4.1. Configuraciones experimentales (OE3).} La misma aplicación se evalúa
en cuatro configuraciones que solo difieren en la capa de defensa:

\begin{itemize}
        \item \textbf{Configuración 1 (control):} aplicación sin protección (con
    \texttt{LOGSGUARDIAN\_DISABLED=true}), para confirmar que cada payload es explotable y
    establecer la línea base.
    \item \textbf{Configuración 2:} aplicación con \textit{logSguarDian} en modo \texttt{block}.
    \item \textbf{Configuración 3a:} aplicación detrás de Nginx con ModSecurity v3 \cite{owasp2026modsecurity} y OWASP CRS
    \cite{owaspcrs} en nivel de paranoia 1 y modo de bloqueo, sin \textit{logSguarDian}.
    \item \textbf{Configuración 3b:} ambas capas superpuestas. El registro de auditoría de
    ModSecurity y el registro de eventos de \textit{logSguarDian} se cruzan para atribuir cada
    bloqueo a la capa que lo produjo, dado que ambas responden con un código de bloqueo.
\end{itemize}

Las cuatro configuraciones se levantan con el mismo árbol de la aplicación, el mismo
\texttt{Dockerfile} y el mismo \texttt{docker-compose.yml}. La librería se desactiva en las
Configuraciones 1 y 3a con la variable de entorno \texttt{LOGSGUARDIAN\_DISABLED=true}, sin
modificar ningún archivo con seguimiento. Las Configuraciones 1 y 2 publican la aplicación
directamente con \texttt{docker-compose.direct.yml}; en la 3a y la 3b el WAF se coloca delante, y
un archivo adicional (\texttt{docker-compose.waf-audit-override.yml}) solo redirige el volumen del
registro de auditoría. El WAF es la imagen \texttt{owasp/modsecurity-crs:nginx-alpine}, con
Nginx 1.30.5, ModSecurity 3.0.17 y OWASP CRS 4.30.0, nivel de paranoia 1,
\texttt{SecRuleEngine On}, umbral de anomalía de entrada de 5 y sin exclusiones activas. La
etiqueta de esa imagen es flotante, por lo que se cita su digest
(\texttt{sha256:f12da37e...}) y no la etiqueta. En las Configuraciones 2 y 3b se verificó, por
separado, que el paquete \texttt{logsguardian} 0.1.1 es el del registro de npm (SHA-256 del
archivo \texttt{abf8d5c4...}) y que los modelos son los del par~C (\texttt{25b407e6...} y
\texttt{5be28ee8...}). Una medición anterior de la Configuración~3a, hecha el 10 de octubre de 2026
en otra rama de la aplicación y con el ejecutor anterior a la corrección, se descartó y se
reemplazó por la de este procedimiento.

Se eligió el nivel de paranoia 1 porque es la configuración realista para la población
objetivo: organizaciones sin equipo de seguridad que pueda ajustar los falsos positivos de los
niveles superiores.

\textbf{Procedimiento 4.2. Auditoría de fuga entre evaluación y entrenamiento.}
\label{subsec:met-fuga}
La fuga de datos, es decir, la presencia en el conjunto de evaluación de información disponible
durante el entrenamiento, es una causa frecuente de sobreestimación del desempeño en estudios
basados en aprendizaje automático \cite{kapoor2023}. Un script versionado
(\texttt{leakage\_check.py}) cruza cada corpus de evaluación contra cada partición en dos
niveles: coincidencia exacta (hash del payload) y coincidencia de plantilla. La plantilla
normaliza la decodificación URL, las mayúsculas y los dígitos; para Command Injection reemplaza
además el marcador aleatorio de SecLists por un símbolo fijo. La auditoría confirmó que el corpus
de la Ronda~4 tenía fuga en Command Injection (los 200 payloads coinciden por plantilla con la
fuente SecLists usada en entrenamiento) y que el conjunto de 500 payloads de la suite E2E se
había muestreado del corpus completo (entre 70 y 75 coincidencias exactas con entrenamiento por
cada 100 payloads de cada clase). En consecuencia, la detección de payloads no vistos se mide
exclusivamente con el corpus limpio M2; el conjunto E2E se usa para verificar el funcionamiento
extremo a extremo de la librería.

\textbf{Procedimiento 4.3. Construcción del corpus limpio (OE3).} El corpus M2 se genera con
scripts versionados y semilla fija, de modo que dos ejecuciones sobre las mismas entradas
producen el mismo corpus: SQLi, Path Traversal y Command Injection con
\texttt{build\_corpus.py} y la semilla 20260923, y XSS con \texttt{generate\_xss.py} y la semilla
20261025, que difiere de la anterior en la semilla y en la fecha (25 de octubre de 2026):

\begin{itemize}
    \item \textbf{SQLi:} sqlmap \cite{sqlmap} v1.10.9, con los parámetros \texttt{--batch
    --level=3 --risk=2 --technique=BEUSTQ}, contra los dos endpoints SQLi. Se registra cada
    petición enviada mediante su archivo de tráfico; se usan las formas de payload que sqlmap
    intentó, no su veredicto de explotabilidad, porque los bloqueos 403 del middleware alteran
    sus heurísticas.
    \item \textbf{Path Traversal/LFI:} generador propio que combina profundidades de 2 a 12
    niveles, ocho estilos de codificación (crudo, barra invertida, codificación URL simple, doble
    y de mayúsculas mixtas, UTF-8 sobrelargo, evasión \texttt{....//} y sufijo de byte nulo) y 28
    archivos objetivo de Linux, \texttt{/proc}, la propia aplicación y Windows.
    \item \textbf{Command Injection:} generador propio, deliberadamente no derivado de SecLists,
    que combina 12 separadores de shell y 4 envolturas de sustitución con 20 comandos de
    post-explotación sobre el sumidero \texttt{ping -c 3 \$\{host\}}.
        \item \textbf{XSS:} generador propio determinista (\texttt{generate\_xss.py}, semilla 20261025),
    no derivado de listas públicas, que combina cuatro dimensiones: 8 contextos de inyección, 9
    vectores de ejecución, 7 funciones ejecutadas y 10 codificaciones o evasiones, con un mapeo
    explícito de los vectores válidos en cada contexto. Genera 400 candidatos, de los que
    \texttt{sample\_xss.py} toma 200 por muestreo estratificado. Una quinta dimensión, que se aplica
    al muestrear, combina 3 transportes (\texttt{application/x-www-form-urlencoded}, JSON y
    \texttt{multipart/form-data}) con 3 campos de destino (\texttt{content} y \texttt{title} de
    \texttt{POST /posts}, y \texttt{bio} de \texttt{POST /profile}). Ningún valor de una dimensión
    supera el 30~\% del corpus, salvo el transporte y el campo, que solo tienen tres valores y para
    los que el mínimo posible del valor más frecuente es 33.3~\%; en esas dos dimensiones se fijó un
    tope de 34~\%. El protocolo se congeló antes de la medición (\texttt{docs/protocol-xss.md}). El
    payload único que confirmó OWASP ZAP \cite{owaspzap} con un \textit{Active Scan} dirigido por su
    API REST contra \texttt{POST /posts} (campos \texttt{content} y \texttt{title}) y
    \texttt{POST /profile} (campo \texttt{bio}), con sesión autenticada, se conserva como hallazgo de
    vulnerabilidad real y no como corpus de medición: ZAP reporta una alerta por cada combinación de
    parámetro y técnica confirmada como vulnerable, no cada payload probado.
\end{itemize}

Todo candidato de SQLi, Path Traversal y Command Injection pasa por un filtro de exclusión antes de entrar al corpus: se descarta si su hash
(o el de su forma codificada o decodificada) coincide con algún payload de entrenamiento de la
misma clase o, en Command Injection, si su plantilla coincide con alguna plantilla de la fuente
SecLists o de la clase \texttt{cmdi} de entrenamiento. Los payloads coincidentes se descartan, no
se marcan, y cada descarte queda registrado con su motivo y la fuente coincidente. De cada
categoría se seleccionan 200 payloads finales; el corpus resultante tiene 800 peticiones. En XSS,
el filtro de exclusión (\texttt{exclusion\_filter\_xss.py}) se aplicó sobre los 200 payloads ya
muestreados, es decir, después y no antes del muestreo, y descartó 0. La auditoría posterior de
fuga (\texttt{leakage\_check\_xss.py}), que mide contención de subcadena normalizada, encontró 3
payloads coincidentes que el filtro no había atrapado, y por regla del protocolo no se corrigieron.

\textbf{Procedimiento 4.4. Medición de cobertura (OE3).} Un ejecutor (\texttt{run\_large\_corpus.py})
envía cada payload del corpus a su endpoint objetivo con la sesión correspondiente y registra, por
petición, el código de respuesta y, cuando aplica, la clase y la confianza del evento almacenado.
Una petición cuenta como detectada solo si el middleware la bloquea con HTTP 403. Las que el
Isolation Forest marca como anómalas sin bloquearlas (\texttt{pass\_anomaly}), las que terminan por
tiempo de espera y las que devuelven HTTP 500 cuentan como no detectadas. La medición principal usa
el paquete 0.1.0 (par~A), con el ejecutor en la versión anterior a la corrección del 10 de octubre
de 2026; esa corrección evita volver a codificar los payloads de Path Traversal y de Command
Injection que ya llegan codificados, y desde entonces el ejecutor envía sus bytes tal cual. Como
referencia, se repitió sobre el mismo corpus con el tarball de
validación 22550ea (par~C), verificando dentro del contenedor el hash de los modelos cargados. En la
Ronda~5 original (9 y 10 de octubre de 2026), las cuatro configuraciones del Procedimiento~4.1 se
ejecutaron sobre el corpus limpio de 601 peticiones; la Configuración~3a se midió entonces en una
rama distinta de la aplicación, y esa medición se descartó (Procedimiento~4.1). La Configuración~2
se repitió además con el paquete 0.1.1, instalado desde el registro de npm, y con el ejecutor
corregido. El 25 de octubre de 2026 las cuatro configuraciones se midieron de nuevo sobre el mismo
árbol de la aplicación, con el ejecutor corregido y con el corpus combinado de 800 peticiones (las
600 de SQLi, Path Traversal y Command Injection, idénticas una a una a las del corpus original, más
200 de XSS), en tres corridas por configuración con el contenedor y la base de datos reiniciados
entre una y otra. Los payloads de XSS se envían con un emisor propio
(\texttt{send\_xss\_corpus.py}) en los tres transportes, que aplica la convención del ejecutor
corregido a los payloads ya codificados.

\textbf{Procedimiento 4.5. Benchmark de rendimiento (objetivo general).} El benchmark se ejecuta
en Linux nativo mediante un flujo de GitHub Actions (\texttt{latency-benchmark.yml}) en un
ejecutor \texttt{ubuntu-latest}, que documenta el hardware del ejecutor en cada corrida. La
aplicación de referencia se levanta con Docker Compose, y la condición se alterna entre línea
base y middleware activo mediante la variable \texttt{LOGSGUARDIAN\_DISABLED}, recreando el
contenedor, sin modificar ningún otro parámetro. Se usa Artillery 2.0.34 \cite{artillerydocs}
con tres escenarios:

\begin{itemize}
    \item \textbf{Normal:} navegación legítima (inicio de sesión, listado, detalle y perfil) a
    25 usuarios virtuales por segundo, que con cuatro peticiones por usuario produce
    $\approx$100 peticiones/s.
    \item \textbf{Ataque:} tráfico mixto concurrente, 70~\% de navegación legítima y 30~\% de
    escenarios con payloads de las cuatro categorías, a 30 usuarios virtuales por segundo
    ($\approx$100 peticiones/s).
    \item \textbf{Volumen:} el escenario normal sobre una base de datos sembrada con 2{,}010
    registros, para representar una aplicación con datos reales. Este escenario no estaba en el
    protocolo; se agregó para contrastar la hipótesis de que una línea base más lenta reduce el
    impacto relativo del middleware, y sus resultados se comparan con el mismo criterio de 15~ms.
\end{itemize}

Cada corrida tiene una fase de calentamiento de 15~s a 5 usuarios/s, descartada del análisis
porque los hilos del Isolation Forest tardan cerca de 5~s en quedar listos, y una fase medida de
60~s. Cada combinación de escenario y condición se repite cinco veces. Durante toda la ejecución,
un proceso paralelo muestrea cada 5~s la memoria y el uso de CPU del contenedor de la aplicación
mediante \texttt{docker stats}. Un script posterior (\texttt{analyze\_resources.py}) asigna cada
muestra a su repetición con las marcas de tiempo de inicio y fin de cada reporte de Artillery,
descarta los primeros 15~s de cada repetición y calcula, para cada par de repeticiones
equivalentes, la diferencia entre la condición con middleware y la línea base. La evidencia
cruda de la corrida (reportes de Artillery, muestras de recursos y hashes de los modelos) está
archivada en el repositorio de la aplicación de referencia. El procedimiento de agregación se
detalla en la Sección~\ref{sec:met-datos}.

\textbf{Procedimiento 4.6. Estudio de falsos positivos con navegadores reales (OE3).} El
procedimiento surgió de un falso positivo reportado al navegar con un navegador real hacia
\texttt{GET /login}. Las corridas de Artillery usan una única cadena \texttt{User-Agent} fija y no
cubren ese caso. Un script (\texttt{ua\_fp\_study.py}) envía peticiones sin consulta ni cuerpo a
\texttt{/login}, \texttt{/posts} y \texttt{/profile} con cadenas \texttt{User-Agent} de Chrome,
Firefox y Safari en macOS y Windows, con y sin sesión, y registra el código de respuesta de cada
combinación.

\textbf{Procedimiento 4.7. Validación de generalización en aplicaciones de terceros (objetivo
general).} Sustituye a la validación con startups guatemaltecas prevista en el protocolo y se
realizó por recomendación del asesor. Se evaluaron cuatro aplicaciones Node.js deliberadamente vulnerables que no participaron en la
calibración ni en el entrenamiento: DVNA, node-api-goat, dvws-node y OWASP Juice Shop. En las tres
primeras se integró el paquete publicado en npm (\texttt{npm install logsguardian}, entonces 0.1.0);
OWASP Juice Shop la evaluó el componente de IA con una copia local del paquete que contiene el
par~C. Las cuatro difieren en arquitectura (renderizado de
plantillas, API pura sin base de datos, SOAP/GraphQL y aplicación de página única). En las tres primeras se confirma primero la explotabilidad sin protección y luego se mide, con el
middleware en modo \texttt{block} y configuración por defecto, la tasa de bloqueo sobre ataques y la tasa de bloqueo
sobre tráfico legítimo representativo de la aplicación. En Juice Shop, en modo \texttt{monitor},
se aisló el efecto de la cabecera \texttt{User-Agent} variando solo esa cabecera (cURL frente a
una cadena de navegador) sobre el mismo conjunto de peticiones.

\subsection{Fase 5: aplicación de feedback}
\label{subsec:met-f5}
Cada hallazgo de la Fase~4 que afecta a una métrica se convierte en una tarea del tablero, con
su propia rama y pull request, y sigue un ciclo fijo: reproducción con evidencia antes de
cualquier cambio; corrección; nueva medición de la métrica afectada; y registro del valor antes
y después en la bitácora de decisiones. Si una corrección requiere reentrenar el modelo, el
conjunto de prueba se bloquea de nuevo antes del reentrenamiento (Procedimiento~1.5). Una
corrección no puede justificarse con la misma lectura de prueba que la motivó.

\subsection{Fase 6: documentación}
\label{subsec:met-f6}
Se consolida un documento de cumplimiento por objetivo, en el que cada afirmación remite al
artefacto que la respalda (script, archivo de resultados, ejecución de CI o captura) y cada
cifra histórica superada se marca como tal. El informe final, el README publicado en npm y la
documentación técnica se derivan de ese documento.

% ------------------------------------------------------------------------------
\section{Técnicas e instrumentos}
\label{sec:met-instrumentos}

El Cuadro~\ref{tab:met-instrumentos} relaciona cada técnica de obtención de datos con el
instrumento que la implementa y con el dato que produce.

\begin{longtable}{|L{2.7cm}|L{3.7cm}|L{3.6cm}|L{2.6cm}|}
\hline
\textbf{Técnica} & \textbf{Instrumento} & \textbf{Dato producido} & \textbf{Procedimiento} \\ \hline
\endfirsthead
\multicolumn{4}{l}{\footnotesize\textit{(continuación)}} \\ \hline
\textbf{Técnica} & \textbf{Instrumento} & \textbf{Dato producido} & \textbf{Procedimiento} \\ \hline
\endhead
\multicolumn{4}{r}{\footnotesize\textit{Continúa en la página siguiente}} \\
\endfoot
\endlastfoot
Análisis documental & Páginas oficiales OWASP, CWE y MITRE ATT\&CK; registros de licencia
de cada fuente & Mapeo de vectores; tabla de auditoría de fuentes & 1.1, 1.2 \\ \hline
Transformación determinista & Analizadores por fuente, \texttt{label\_map.yaml},
\texttt{unify.py}, \texttt{split.py} & Corpus etiquetado, particiones y hash de bloqueo & 1.3--1.5 \\ \hline
Extracción de características & CLI de \texttt{packages/extractor} (TypeScript) & Matriz de
características por petición & 1.3, F2 \\ \hline
Evaluación sobre conjunto bloqueado & \texttt{evaluate\_test.py}, scikit-learn
\cite{pedregosa2011}, \texttt{onnxruntime} & Precisión, recall, F1, AUC-ROC; métricas del
Isolation Forest; manifiesto de sumas & 2.2 \\ \hline
Barrido de umbral & Notebooks de calibración sobre validación & Curvas precisión/recall/FP por
umbral & 2.1 \\ \hline
Pruebas automatizadas & Jest, Supertest, GitHub Actions & Resultado de pruebas por commit & 3.2 \\ \hline
Auditoría de fuga & \texttt{leakage\_check.py} & Coincidencias exactas y por plantilla por
clase y partición & 4.2 \\ \hline
Pruebas de penetración controlada & sqlmap \cite{sqlmap}, generadores propios, filtro de
exclusión y ejecutor del corpus & Corpus limpio; veredicto por
payload y configuración & 4.3, 4.4 \\ \hline
WAF de comparación & ModSecurity v3 + OWASP CRS PL1 \cite{owaspcrs} & Registro de auditoría
del WAF & 4.1 \\ \hline
Prueba de carga & Artillery 2.0.34, script de orquestación de la suite y flujo de CI del
benchmark & Latencia p95 por repetición; errores y tiempos de espera
excedidos & 4.5 \\ \hline
Monitoreo de recursos & \texttt{docker stats} y script de análisis de memoria & Serie temporal de memoria y CPU del contenedor & 4.5 \\ \hline
Experimento de cabecera & Script de estudio de \texttt{User-Agent}; corridas cURL frente a
navegador en Juice Shop & Veredicto por cadena \texttt{User-Agent} y ruta & 4.6, 4.7 \\ \hline
Registro de eventos & SQLite del middleware; CLI \texttt{attacks} y \texttt{endpoints} & Clase,
confianza y veredicto por petición & 4.4, 4.7 \\ \hline
\caption[Técnicas e instrumentos de obtención de datos]{Técnicas e instrumentos de obtención de datos, dato producido y procedimiento asociado}
\label{tab:met-instrumentos}
\end{longtable}
{\footnotesize Nota. Elaboración propia.}

OWASP ZAP \cite{owaspzap}, previsto en el protocolo para generar payloads XSS, está instalado como
aplicación gráfica; no se automatizó mediante el navegador, sino a través de la API REST que expone
su proxy, que permitió crear el contexto, autenticar al usuario y lanzar el \textit{Active Scan}
sin intervención manual. Su resultado es un único payload confirmado, que se conserva como hallazgo de explotabilidad
(n=1; Sección~\ref{sec:met-restricciones}) y no como corpus de medición: la cobertura de XSS se mide
con el generador propio de 200 payloads (Procedimiento~4.3). La herramienta
clinic.js, prevista para medir CPU, se sustituyó por \texttt{docker stats}: mide el contenedor
completo, puede ejecutarse sin intervención dentro del flujo de CI y produce una serie temporal
alineada con las corridas de Artillery.

% ------------------------------------------------------------------------------
\section{Recursos}
\label{sec:met-recursos}

\subsection{Recursos humanos}
El componente de ciberseguridad estuvo a cargo de un único integrante, con dedicación estimada
de 20 horas semanales en las Fases 1 a 3 y de 15 horas semanales en las Fases 4 a 6. El
componente de IA estuvo a cargo del segundo integrante, que además actuó como revisor obligatorio
de los pull requests. La asesoría académica estuvo a cargo del asesor asignado, cuyas revisiones
del documento de cumplimiento motivaron varias de las correcciones descritas en la
Sección~\ref{sec:met-cambios}. No participó personal externo.

\subsection{Recursos tecnológicos}
El Cuadro~\ref{tab:met-recursos} actualiza los recursos tecnológicos del protocolo con las
herramientas que se usaron en la práctica.

\begin{longtable}{|L{3.0cm}|L{4.5cm}|L{5.3cm}|}
\hline
\textbf{Categoría} & \textbf{Herramienta} & \textbf{Función en el proyecto} \\ \hline
\endfirsthead
\multicolumn{3}{l}{\footnotesize\textit{(continuación)}} \\ \hline
\textbf{Categoría} & \textbf{Herramienta} & \textbf{Función en el proyecto} \\ \hline
\endhead
\multicolumn{3}{r}{\footnotesize\textit{Continúa en la página siguiente}} \\
\endfoot
\endlastfoot
Runtime y framework & Node.js, Express & Plataforma de la librería y de las aplicaciones de prueba \\ \hline
Monorepo & pnpm workspaces, TypeScript & Paquetes \texttt{extractor}, \texttt{core} y \texttt{mlops} \\ \hline
Inferencia & \texttt{onnxruntime-node}, \texttt{worker\_threads} & Ejecución de los modelos sin bloquear el Event Loop \\ \hline
Persistencia & SQLite (\texttt{better-sqlite3}) & Registro de eventos y de webhooks \\ \hline
Entorno de prueba & Docker Compose, PostgreSQL & Aplicación de referencia reproducible \\ \hline
Generación de tráfico & Artillery 2.0.34 & Carga legítima y mixta \\ \hline
Simulación de ataques & sqlmap 1.10.9; generadores propios en Python & Payloads de SQLi, Path Traversal y Command Injection \\ \hline
WAF de comparación & Nginx + ModSecurity v3 + OWASP CRS & Configuraciones 3a y 3b \\ \hline
Análisis de datos y ML & Python, pandas, scikit-learn, Jupyter & Partición, evaluación, calibración y auditoría de fuga \\ \hline
Pruebas & Jest, Supertest & Pruebas unitarias, de integración y E2E \\ \hline
Integración continua & GitHub Actions (\texttt{ubuntu-latest}) & Pruebas por pull request, publicación en npm y benchmark de rendimiento en Linux nativo \\ \hline
Gestión & Notion (tablero Kanban), Git, GitHub & Gestión de tareas, versionado y revisión cruzada \\ \hline
Distribución & Registro público de npm & Publicación de \texttt{logsguardian} \\ \hline
\caption[Recursos tecnológicos utilizados]{Recursos tecnológicos utilizados en el proyecto}
\label{tab:met-recursos}
\end{longtable}
{\footnotesize Nota. Elaboración propia. Todas las herramientas son de código abierto o de uso
gratuito para el proyecto.}

\subsection{Recursos externos y financieros}
Los recursos externos son las fuentes públicas del corpus, las cuatro aplicaciones vulnerables
de terceros, el registro de npm y los ejecutores de GitHub Actions. El componente de
ciberseguridad no requirió licencias de pago. La infraestructura en la nube contemplada en el
protocolo para la validación con startups no se utilizó, porque esa validación no se realizó.

% ------------------------------------------------------------------------------
\section{Recolección, procesamiento y validación de datos}
\label{sec:met-datos}

El Cuadro~\ref{tab:met-datos} describe, para cada flujo de datos, cómo se recolecta, cómo se
procesa y qué verificación debe superar antes de considerarse evidencia válida para el capítulo
de Resultados.

\begin{longtable}{|L{2.3cm}|L{3.4cm}|L{3.4cm}|L{3.7cm}|}
\hline
\textbf{Flujo} & \textbf{Recolección} & \textbf{Procesamiento} & \textbf{Validación} \\ \hline
\endfirsthead
\multicolumn{4}{l}{\footnotesize\textit{(continuación)}} \\ \hline
\textbf{Flujo} & \textbf{Recolección} & \textbf{Procesamiento} & \textbf{Validación} \\ \hline
\endhead
\multicolumn{4}{r}{\footnotesize\textit{Continúa en la página siguiente}} \\
\endfoot
\endlastfoot
Corpus etiquetado (M1) & Descarga de las fuentes públicas (Cuadro~\ref{tab:met-origenes}) & Esquema canónico,
etiquetado por mapeo, deduplicación exacta, extracción de características, partición
estratificada & Cero solapamiento entre particiones; tolerancias de clase ($\pm$2~\%) y fuente
($\pm$5~\%); hash de prueba versionado antes de entrenar \\ \hline
Métricas de clasificación & Lectura del conjunto de prueba con los binarios ONNX publicados
& Reporte de clasificación y AUC-ROC uno contra el resto & Hash de prueba coincidente;
manifiesto de sumas de \texttt{test.parquet}, \texttt{rf.onnx} e \texttt{if.onnx}; umbrales
tomados de los metadatos del modelo \\ \hline
Paridad ONNX & Predicciones de Python y de Node.js sobre el mismo conjunto & Diferencia máxima
absoluta por salida & Prueba automatizada en CI; diferencia $<$0.1~\% \\ \hline
Corpus limpio (M2) & sqlmap y generadores propios con semilla fija & Filtro de exclusión
exacto y por plantilla; selección de 200 por categoría & Sin coincidencias por hash exacto
tras el filtro, por construcción; la auditoría posterior por subcadena normalizada encontró 4
(SQLi), 7 (Path Traversal), 2 (Command Injection) y 3 (XSS) coincidencias de 200, y cinco de las
siete de Path Traversal son también literales; reproducibilidad (dos ejecuciones con igual resultado) \\ \hline
Cobertura & Envío por HTTP real del corpus a la aplicación en cada configuración & Conteo de
bloqueos (HTTP 403) por categoría; anomalías sin bloqueo, tiempos de espera y errores 500 cuentan
como no detectados & Explotabilidad confirmada en Configuración~1; atribución por
capa en 3b mediante el registro de auditoría del WAF \\ \hline
Latencia & Reportes JSON de Artillery por repetición & p95 por repetición sobre la fase medida;
mediana de las cinco repeticiones; $\Delta$p95 absoluto y relativo & Exclusión del
calentamiento; tasa real de peticiones $\approx$100/s verificada en el reporte; conteo de
tiempos de espera excedidos a partir del resumen final de cada repetición \\ \hline
Memoria y CPU & Muestreo cada 5~s de \texttt{docker stats} & Asignación de muestras a cada
repetición por marca de tiempo; media por condición (MiB); incremento pareado con middleware menos línea
base & Ventana total $\geq$30~min; se excluyen calentamiento y muestras fuera de una repetición;
contraste entre la ventana medida y la repetición completa \\ \hline
Falsos positivos con navegador (M4) & Script con cadenas \texttt{User-Agent} reales &
Proporción de combinaciones bloqueadas & Sesión y ruta controladas; una sola variable manipulada \\ \hline
Aplicaciones de terceros (M5) & Integración del paquete publicado y scripts de ataque y de
tráfico legítimo por aplicación & Tasa de bloqueo en ataques y en tráfico legítimo, reportadas
siempre juntas & Explotabilidad previa confirmada; evidencia cruda (scripts y JSON) copiada al
repositorio \\ \hline
\caption[Recolección, procesamiento y validación de datos]{Recolección, procesamiento y validación de cada flujo de datos}
\label{tab:met-datos}
\end{longtable}
{\footnotesize Nota. Elaboración propia.}

Tres decisiones de procesamiento requieren justificación explícita, porque afectan la
interpretación de los resultados:

\begin{itemize}
    \item \textbf{Agregación de las repeticiones de latencia.} De las cinco repeticiones de cada condición se reporta la mediana. Se usa la mediana porque es menos sensible que la media a una repetición atípica
    causada por el entorno de ejecución.
    \item \textbf{Cálculo del p95 por repetición.} Artillery reporta percentiles por intervalos
    de aproximadamente 10~s. El analizador calcula el p95 de cada repetición como el promedio,
    ponderado por número de peticiones, de los p95 de cada intervalo posterior al calentamiento,
    no como el percentil exacto sobre todas las muestras. Es una aproximación y se declara como
    tal; su efecto es mayor en el escenario de volumen, donde la latencia crece durante la corrida.
    \item \textbf{Lectura conjunta de detección y falsos positivos.} En las aplicaciones de
    terceros, una tasa de bloqueo alta en ataques no se reporta sin la tasa de bloqueo en tráfico
    legítimo. Un sistema que bloquea casi todo obtiene un recall alto sin discriminar.
\end{itemize}

% ------------------------------------------------------------------------------
\section{Métricas e indicadores}
\label{sec:met-metricas}

Las métricas de clasificación se calculan por clase. Sean $TP_c$, $FP_c$ y $FN_c$ los
verdaderos positivos, falsos positivos y falsos negativos de la clase $c$ (Ecuación~\eqref{eq:met-clasificacion}):

\begin{equation}
P_c = \frac{TP_c}{TP_c + FP_c}, \qquad
R_c = \frac{TP_c}{TP_c + FN_c}, \qquad
F1_c = \frac{2\,P_c\,R_c}{P_c + R_c}
\label{eq:met-clasificacion}
\end{equation}

El F1 es la métrica primaria de aceptación porque el corpus está desbalanceado: la exactitud
global premiaría a un clasificador que favoreciera a la clase mayoritaria, mientras que el F1
penaliza tanto los ataques no detectados como los bloqueos de tráfico legítimo. El AUC-ROC uno
contra el resto se reporta como medida de discriminación independiente del umbral
\cite{fawcett2006}. Para el rendimiento (Ecuación~\eqref{eq:met-latencia}):

\begin{equation}
\Delta p95_{\mathrm{abs}} = \tilde{p}95_{\mathrm{activo}} - \tilde{p}95_{\mathrm{base}}, \qquad
\Delta p95_{\mathrm{rel}} = \frac{\Delta p95_{\mathrm{abs}}}{\tilde{p}95_{\mathrm{base}}} \times 100\,\%
\label{eq:met-latencia}
\end{equation}

\noindent donde $\tilde{p}95$ es la mediana de las cinco repeticiones de una condición.

Para las proporciones observadas (tasas de bloqueo y de detección) se reportan intervalos de confianza de Wilson al 95~\% \cite{wilson1927probable}, preferidos al intervalo normal (Wald) porque mantienen una cobertura cercana a la nominal con muestras pequeñas o proporciones cercanas a 0 o 1 \cite{brown2001interval}.

El criterio de aceptación de latencia es el de la Sección~\ref{subsec:alcance-tecnico}: un
$\Delta p95_{\mathrm{abs}}$ de 15~ms como máximo, tanto bajo carga normal como bajo carga de
ataque. Se expresa en milisegundos y no como porcentaje de la línea base porque un porcentaje
depende por completo de la latencia propia de la aplicación anfitriona. Con una línea base de 2 a
5~ms, como la de la aplicación de referencia, un margen de 5~\% equivale a una fracción de
milisegundo, menor que la variación entre repeticiones del entorno de medición, de modo que el
criterio dejaría de distinguir un middleware costoso de uno ligero. Además, el mismo middleware
obtendría veredictos opuestos según la aplicación en que se instale: en una aplicación más lenta,
el mismo costo se diluye y el porcentaje baja. Una cifra absoluta mide directamente el costo que
el middleware agrega a cada petición, que es la magnitud que percibe el usuario final y la que un
equipo sin especialista en seguridad puede comparar con su propio presupuesto de latencia.
$\Delta p95_{\mathrm{rel}}$ se reporta como indicador complementario.

El Cuadro~\ref{tab:met-metricas} define todas las métricas, su criterio y el objetivo al que se
vinculan.

\begin{longtable}{|L{2.8cm}|L{4.0cm}|L{2.7cm}|L{2.0cm}|L{1.0cm}|}
\hline
\textbf{Indicador} & \textbf{Definición operacional} & \textbf{Criterio} & \textbf{Instrumento} & \textbf{Obj.} \\ \hline
\endfirsthead
\multicolumn{5}{l}{\footnotesize\textit{(continuación)}} \\ \hline
\textbf{Indicador} & \textbf{Definición operacional} & \textbf{Criterio} & \textbf{Instrumento} & \textbf{Obj.} \\ \hline
\endhead
\multicolumn{5}{r}{\footnotesize\textit{Continúa en la página siguiente}} \\
\endfoot
\endlastfoot
Tamaño del corpus & Filas totales en las tres particiones & $\geq$100{,}000 & Partición (1.5) & OE2 \\ \hline
Representación de clases & Clases con muestras en cada partición & 5/5 en cada partición & Partición (1.5) & OE1, OE2 \\ \hline
Integridad de la partición & Hashes compartidos entre particiones & 0 & Partición (1.5) & OE2 \\ \hline
F1 por categoría & $F1_c$ sobre el conjunto de prueba bloqueado & $\geq$0.80 en al menos 3 de 4 categorías & Lectura de prueba (2.2) & OE3 \\ \hline
Precisión y recall por categoría & $P_c$, $R_c$ sobre el conjunto de prueba & Reportados, subordinados al F1 & Lectura de prueba (2.2) & OE3 \\ \hline
AUC-ROC por categoría & Área bajo la curva ROC uno contra el resto & Referencia, sin umbral & Lectura de prueba (2.2) & OE3 \\ \hline
Detección de anomalías del IF & $TP/(TP+FN)$ sobre ataques de prueba, umbral de los metadatos & $\geq$50~\% & Lectura de prueba (2.2) & OE3 \\ \hline
Falsos positivos del IF & $FP/(FP+TN)$ sobre legítimos de prueba & $\leq$10~\% & Lectura de prueba (2.2) & OE3 \\ \hline
Paridad ONNX & Diferencia máxima entre predicciones Python y Node.js & $<$0.1~\% & Prueba de paridad & OG \\ \hline
Cobertura por categoría & Payloads del corpus limpio bloqueados con HTTP 403 / payloads enviados & $\geq$80~\% por categoría & Ejecutor del corpus (4.4) & OE3 \\ \hline
Contribución en defensa en profundidad & Ataques no bloqueados por el WAF que bloquea \textit{logSguarDian} / ataques no bloqueados por el WAF & Descriptivo & Configuraciones 3a/3b & OG \\ \hline
Falsos positivos con navegador & Combinaciones ruta--\texttt{User-Agent} legítimas bloqueadas / combinaciones probadas & Descriptivo; se reporta junto al F1 & Estudio UA (4.6) & OE3 \\ \hline
Generalización a terceros & Tasa de bloqueo en ataques y en tráfico legítimo por aplicación & Descriptivo; siempre en pareja & Scripts por aplicación & OG \\ \hline
Extracción de características & Tiempo por petición, aislado en su hilo & $\leq$5~ms & Benchmark del extractor & OG \\ \hline
Latencia adicional & $\Delta p95_{\mathrm{abs}}$, escenarios normal y de ataque (criterio) y de volumen (referencia) & $\leq$15~ms; $\Delta p95_{\mathrm{rel}}$ complementario & Artillery, CI & OG \\ \hline
Tiempos de espera excedidos & Peticiones con \texttt{ERR\_SOCKET\_TIMEOUT} / peticiones enviadas & Descriptivo (escenario de volumen) & Artillery, CI & OG \\ \hline
Memoria adicional & Incremento pareado de memoria del contenedor con middleware respecto de la línea base (MiB) & $\leq$150~MB (143.1~MiB) & \texttt{docker stats} & OG \\ \hline
Incremento de CPU & Uso de CPU con middleware frente a la línea base & $\leq$15~\% relativo & \texttt{docker stats} & OG \\ \hline
\caption[Métricas e indicadores de evaluación]{Métricas e indicadores: definición operacional, criterio de aceptación, instrumento y objetivo vinculado}
\label{tab:met-metricas}
\end{longtable}
{\footnotesize Nota. Elaboración propia con base en el protocolo de graduación y en
\cite{fawcett2006}. OG: objetivo general.}

% ------------------------------------------------------------------------------
\section{Matriz de trazabilidad}
\label{sec:met-trazabilidad}

El Cuadro~\ref{tab:met-trazabilidad} relaciona cada objetivo con las actividades que lo
abordan, los datos que estas producen, las métricas que se calculan sobre esos datos y el
resultado esperado, que es lo que el capítulo de Resultados contrasta.

\begin{longtable}{|L{1.5cm}|L{2.6cm}|L{2.7cm}|L{2.6cm}|L{3.3cm}|}
\hline
\textbf{Objetivo} & \textbf{Actividades} & \textbf{Datos} & \textbf{Métricas} & \textbf{Resultado esperado} \\ \hline
\endfirsthead
\multicolumn{5}{l}{\footnotesize\textit{(continuación)}} \\ \hline
\textbf{Objetivo} & \textbf{Actividades} & \textbf{Datos} & \textbf{Métricas} & \textbf{Resultado esperado} \\ \hline
\endhead
\multicolumn{5}{r}{\footnotesize\textit{Continúa en la página siguiente}} \\
\endfoot
\endlastfoot
OE1 & 1.1 Mapeo a OWASP, CWE y MITRE & Cuadro de correspondencia verificado & Vectores mapeados con fuente oficial & Cuatro vectores justificados; cinco clases canónicas que estructuran OE2 y OE3 \\ \hline
OE2 & 1.2--1.5 Auditoría, etiquetado, deduplicación y partición & Corpus M1; reporte de auditoría; hash de bloqueo & Tamaño; representación de clases; integridad de partición & $\geq$100{,}000 muestras etiquetadas con etiquetado determinista y documentado \\ \hline
OE3 (modelo) & 2.1 Calibración; 2.2 Lectura de prueba & Conjunto de prueba bloqueado; manifiesto de sumas & F1, precisión, recall y AUC-ROC por clase; detección y FP del IF & F1 $\geq$0.80 en al menos 3 de 4 categorías; IF $\geq$50~\% de detección con $\leq$10~\% de FP \\ \hline
OE3 (sistema) & 4.1--4.4 Configuraciones, auditoría de fuga, corpus limpio y cobertura & Corpus M2; veredictos por payload y configuración & Cobertura por categoría; contribución en defensa en profundidad & Cobertura $\geq$80~\% por categoría sobre payloads no vistos en entrenamiento \\ \hline
OE3 (falsos positivos) & 4.6 Estudio con navegadores & Muestra M4 & Falsos positivos con navegador & Cuantificación del bloqueo de tráfico legítimo no cubierto por Artillery \\ \hline
OG (rendimiento) & 3.1 Arquitectura en hilos; 4.5 Benchmark & Muestra M3; series de memoria y CPU & $\Delta p95$; tiempos de espera; memoria; CPU; extracción & $\Delta p95_{\mathrm{abs}}\leq$15~ms (normal y ataque; volumen como referencia); $\Delta$ memoria $\leq$150~MB (143.1~MiB); CPU $\leq$15~\% \\ \hline
OG (generalización) & 4.7 Integración en terceros & Muestra M5 & Bloqueo en ataques y en legítimos por aplicación & Evidencia de si la calibración transfiere fuera de la aplicación de referencia \\ \hline
OG (librería) & 3.1--3.2 Desarrollo y verificación; publicación & Pruebas en CI; paquete publicado & Pruebas aprobadas; paridad ONNX & Paquete funcional publicado en npm \\ \hline
\caption[Matriz de trazabilidad metodológica]{Matriz de trazabilidad entre objetivos, actividades, datos, métricas y resultados esperados}
\label{tab:met-trazabilidad}
\end{longtable}
{\footnotesize Nota. Elaboración propia. Los números de actividad remiten a los procedimientos
de la Sección~\ref{sec:met-fases}.}

% ------------------------------------------------------------------------------
\section{Cambios respecto al Marco Metodológico del protocolo}
\label{sec:met-cambios}

El Cuadro~\ref{tab:met-cambios} enumera los cambios respecto al protocolo y la causa de cada
uno.

\begin{longtable}{|L{3.4cm}|L{4.3cm}|L{5.1cm}|}
\hline
\textbf{Elemento} & \textbf{Cambio} & \textbf{Causa} \\ \hline
\endfirsthead
\multicolumn{3}{l}{\footnotesize\textit{(continuación)}} \\ \hline
\textbf{Elemento} & \textbf{Cambio} & \textbf{Causa} \\ \hline
\endhead
\multicolumn{3}{r}{\footnotesize\textit{Continúa en la página siguiente}} \\
\endfoot
\endlastfoot
Criterio de latencia & De $\Delta p95$ relativo (5~\%/10~\%) a absoluto ($\leq$15~ms) & Una cifra en milisegundos mide el costo real por petición; un porcentaje sobre una línea base de 2--5~ms es menor que la variación de medición y cambia según la aplicación anfitriona \\ \hline
Agregación de repeticiones & De media de las tres centrales a mediana & Menor sensibilidad a repeticiones atípicas \\ \hline
Entorno de medición & De Docker Desktop en macOS a Linux nativo en GitHub Actions & La capa de virtualización de Docker Desktop alteraba la latencia medida \\ \hline
Escenario de volumen & Agregado & Contrastar el efecto de una línea base realista sobre el impacto del middleware \\ \hline
Origen del dataset (OE2) & De generado en el entorno Docker a construido con fuentes públicas auditadas, listas públicas de payloads y conjuntos propios; el entorno se usa para evaluar & Las fuentes públicas aportan volumen y diversidad de ataques que el entorno controlado no podía generar en el tiempo disponible \\ \hline
Corpus de evaluación & Corpus de SecLists (Ronda~4) reemplazado por corpus limpio (Ronda~5) & Fuga confirmada en Command Injection respecto al entrenamiento \\ \hline
Política de decisión & El Isolation Forest no bloquea; solo registra anomalías & Una política disyuntiva bloqueaba $\approx$10~\% del tráfico legítimo \\ \hline
Validación externa & Startups guatemaltecas sustituidas por cuatro aplicaciones vulnerables de terceros & No se concretó un acuerdo con una startup en el período de pruebas \\ \hline
Estudio de falsos positivos con navegador & Agregado & Falso positivo observado con un navegador real, no cubierto por Artillery \\ \hline
Instrumento de CPU & clinic.js sustituido por \texttt{docker stats} & Ejecución sin intervención en CI y serie temporal alineada con la carga \\ \hline
Gestión & Kanban en Notion con cuatro columnas, sin la columna Bloqueado; complementado con protección de ramas e integración continua & Simplificar el tablero; convertir la revisión cruzada y las pruebas en controles obligatorios \\ \hline
Reentrenamiento offline (\texttt{mlops}) & Agregado, en el componente de IA & Recomendación del asesor \\ \hline
\caption[Cambios respecto al Marco Metodológico del protocolo]{Cambios respecto al Marco Metodológico del protocolo y su causa}
\label{tab:met-cambios}
\end{longtable}
{\footnotesize Nota. Elaboración propia.}

% ------------------------------------------------------------------------------
\section{Amenazas a la validez}
\label{sec:met-validez}

Las condiciones siguientes limitan lo que los datos pueden demostrar. Se declaran aquí para que
el capítulo de Resultados las tenga en cuenta al interpretar cada cifra:

\begin{itemize}
    \item \textbf{Rebloqueos y lecturas repetidas del conjunto de prueba.} El conjunto de prueba
    se volvió a bloquear en cada generación de modelo (el valor del bloqueo cambió nueve veces, entre
    el 11 de junio y el 23 de septiembre de 2026), y en la generación
    \texttt{rf\_v11} se leyó tres veces, con un cambio del umbral del Isolation Forest entre la
    primera y la última lectura. En \texttt{if\_v2} e \texttt{if\_v5}, los objetivos internos de
    calibración del Isolation Forest se fijaron después de observar resultados de prueba. Ver
    resultados de prueba antes de decidir filtra información hacia el diseño. El manifiesto de
    sumas de verificación solo protege las lecturas desde \texttt{rf\_v11}.
    \item \textbf{Procedencia de los datos restaurados.} Los datos de entrenamiento se
    restauraron desde una copia de respaldo el 23 de septiembre de 2026 sin que el repositorio
    registre su origen. La partición vigente suma 383{,}202 filas, 119 más que el archivo
    unificado (383{,}083); la diferencia coincide con las 119 filas de telemetría curada del
    pipeline de reentrenamiento, que el script de unificación excluye y el pipeline incorpora por
    separado (Cuadro~\ref{tab:met-origenes}), por lo que no queda sin explicar. El 17.8~\%
    de las filas de Command Injection del conjunto de prueba (261 de 1{,}468) proviene de
    conjuntos sintéticos cuyo script generador no está en el repositorio; esa porción del corpus no
    es reproducible con los scripts versionados.
    \item \textbf{Métricas sin umbral de bloqueo.} Las métricas del Random Forest sobre el
    conjunto de prueba asignan la clase de mayor probabilidad, mientras que el middleware bloquea
    solo si esa probabilidad alcanza 0.35. Ambas decisiones coinciden en la mayoría de los casos,
    pero no son idénticas.
    \item \textbf{Calibración en una sola aplicación.} El umbral se calibró sobre la partición
    de validación y se confirmó en la aplicación de referencia. Las aplicaciones de terceros son
    la única evidencia de cómo transfiere a otras formas de tráfico.
    \item \textbf{Corpus de ataque de las aplicaciones de terceros.} Las integraciones en DVNA,
    node-api-goat y dvws-node, realizadas por recomendación del asesor para probar la librería
    fuera de la aplicación de referencia, usaron el conjunto de 500 payloads de la suite E2E, que
    comparte payloads con el entrenamiento (Procedimiento~4.2). Su tasa de bloqueo en ataques
    confirma que la librería instalada desde npm funciona en aplicaciones ajenas al proyecto, pero
    no mide la detección de payloads nuevos; esa medición corresponde al corpus limpio (M2). Sus
    tasas de bloqueo sobre tráfico legítimo no dependen de ese corpus y son la evidencia principal
    de generalización.
    \item \textbf{Sesgo de representación de la cabecera \texttt{User-Agent}.} Solo el 0.4~\%
    de las filas legítimas del corpus tiene un \texttt{User-Agent} de más de 40 caracteres,
    frente al 33.3--97.7~\% de las clases de ataque. Las métricas sobre el conjunto de prueba
    heredan ese sesgo y no reflejan el tráfico de navegadores reales.
    \item \textbf{Generaciones de modelo distintas entre evaluaciones.} Las métricas sobre el
    conjunto de prueba corresponden al par~C, que es el contenido del paquete 0.1.1; el paquete 0.1.0, la
    medición principal de cobertura y el estudio de \texttt{User-Agent}, al par~A; las aplicaciones DVNA, node-api-goat y
    dvws-node, con atribución probable al par~A, porque el paquete y el modelo exactos de cada
    medición no constan en el repositorio; el benchmark de rendimiento y la repetición de la
    Configuración~2 del 9 de septiembre, al par~B; la evaluación en OWASP Juice Shop y las Configuraciones 2 y 3b de la medición unificada del 25 de
    octubre de 2026, al par~C (Cuadro~\ref{tab:pares-modelos}). Los pares no tienen la misma entrada: el par~A consume 67 y 61
    características, y los pares B y C, 69 y 63. Comparten, en cambio, la arquitectura de ejecución
    (un hilo de Random Forest, dos de Isolation Forest, comunicación entre hilos y escritura en
    SQLite), por lo que los resultados de rendimiento describen principalmente esa arquitectura y no
    los pesos de un modelo en particular, aunque no se midieron con el paquete 0.1.1. El paquete 0.1.0
    usa el umbral del Isolation Forest calibrado para \texttt{if\_v9} (0.002486), fijado como
    constante en el código; el 0.1.1 lo lee del archivo \texttt{parity\_report.json} del modelo (0.004206).
    \item \textbf{Evaluación no ciega del par~C en la Ronda~5.} Los resultados de la Ronda~5 con el
    par~A se conocieron el 24 de septiembre de 2026 y el par~C se reconstruyó el 26 de septiembre.
    El historial no vincula ambos hechos, pero tampoco permite descartar que esos resultados
    influyeran en el reentrenamiento, por lo que la cobertura del par~C no constituye una
    evaluación estrictamente ciega; el paquete 0.1.1 contiene el mismo par.
    \item \textbf{Dos defectos de recodificación en la cadena de medición.} Hubo dos defectos
    independientes. El primero, en el middleware, que reconstruía la cadena de consulta a partir de
    \texttt{req.query}, ya decodificada una vez por Express, y le añadía una capa de codificación;
    se corrigió el 10 de octubre de 2026 y la corrección está en el paquete 0.1.1. El segundo, en el
    ejecutor de pruebas (\texttt{run\_large\_corpus.py}), que volvía a codificar los payloads de
    Path Traversal y de Command Injection que los generadores ya entregaban codificados (por
    ejemplo \texttt{\%252e\%252e\%252f} o un separador \texttt{\%0a}), de modo que el payload enviado
    no era el que el generador pretendía; se corrigió el 10 de octubre, limitado a esas dos
    categorías. Ninguna corrección bastó por sí sola (Cuadro~\ref{tab:res-pt-causas}). Las mediciones de
    Path Traversal y de Command Injection de la Ronda~5 con el paquete 0.1.0 y con el tarball de
    validación, y las de la Ronda~4 que usaron el mismo ejecutor, se hicieron con el ejecutor
    anterior. La Configuración~2 con el paquete 0.1.1 y las cuatro configuraciones de la medición
    unificada del 25 de octubre de 2026 se hicieron con el ejecutor corregido. No se midió el
    paquete 0.1.0 con el ejecutor corregido, por lo que no puede separarse cuánto de la mejora de
    85.0~\% a 100.0~\% se debe al modelo, al middleware y al ejecutor. Para el WAF sí se aisló el
    efecto del ejecutor: el ejecutor anterior, ejecutado sobre el árbol y la pila actuales,
    reprodujo el 68.0~\% de Command Injection de la medición del 10 de octubre (136 de 200, con el
    mismo veredicto en los 600 payloads de esa medición), y el ejecutor corregido dio 164 de 200
    (82.0~\%); la diferencia se debe a que el ejecutor anterior volvía a codificar los separadores
    ya codificados, y no a la versión de la aplicación.
        \item \textbf{Corpus de XSS construido después de conocer los resultados de la Ronda~5.} El
    corpus propio de XSS se diseñó el 25 de octubre de 2026, con el conocimiento de que la Ronda~5
    original medía XSS con n=1 y de que las otras categorías ya tenían generadores propios. La
    medición no es ciega, y el 200 de 200 no puede leerse como una evaluación independiente.
    \item \textbf{Sesgo del autor en el generador de XSS.} Las cuatro dimensiones, sus valores y el
    mapeo de los vectores válidos en cada contexto los diseñó el autor del generador, no se
    derivaron de una taxonomía publicada; la cobertura medida depende de las combinaciones que se
    consideraron válidas.
    \item \textbf{Coincidencias con el entrenamiento que el filtro no eliminó.} La auditoría de
    subcadena normalizada (\texttt{leakage\_check.py}, sobre las particiones de entrenamiento,
    validación y prueba) encontró 3 de los 200 payloads de XSS (\texttt{xss\_118},
    \texttt{xss\_254} y \texttt{xss\_342}) con coincidencia, que el filtro, con 0 descartes, no
    atrapó porque solo descarta la igualdad exacta de la forma normalizada del campo completo; por
    regla del protocolo no se corrigieron ni se retiraron. Los tres se bloquean en las
    Configuraciones 2, 3a y 3b, y sin ellos la librería sigue bloqueando 197 de 197.
    \item \textbf{Auditoría de subcadena en las cuatro categorías.} Aplicada el 25 de octubre de
    2026 a los cuatro corpus, la misma auditoría encontró coincidencias normalizadas por subcadena
    en 4 de los 200 payloads de SQLi, 7 de Path Traversal, 2 de Command Injection y 3 de XSS (por
    igualdad del campo completo, 4 de SQLi y ninguno en las demás). Es un criterio permisivo, que
    sirve como límite superior y no como medida de solapamiento, y el filtro de exclusión de la
    Ronda~5 no lo aplicaba.
    \item \textbf{Campos de XSS sin vulnerabilidad confirmada.} De los 200 payloads de XSS, 136 (68
    en \texttt{content} y 68 en \texttt{title}) se dirigen a campos de \texttt{POST /posts} donde
    ZAP no confirmó vulnerabilidad, y sin protección la aplicación los rechaza con HTTP 400; miden
    la detección y no la explotabilidad. Solo los 64 dirigidos a \texttt{bio} apuntan a un campo
    con vulnerabilidad confirmada.
    \item \textbf{Congelación del protocolo de XSS sin historial de git.} Que el protocolo y los
    artefactos del corpus de XSS no cambiaron entre su congelación y la medición se verificó con un
    manifiesto de sumas SHA-256 (\texttt{docs/freeze-xss-20261025.json}), recalculado antes y
    después, y no con el historial de git, porque la tarea se ejecutó sin commits; la garantía no
    es verificable por un tercero a partir de un commit.
    \item \textbf{Atribución por capa en la Configuración~3b por categoría.} El registro de eventos
    de la librería no guarda un identificador de petición, por lo que los bloqueos de la 3b se
    atribuyeron a una de las dos capas por método y ruta, es decir, por categoría y no por payload.
    La atribución al WAF coincide, categoría por categoría, con lo medido por separado en la 3a,
    pero no es un cruce determinista.
    \item \textbf{La librería no se comporta igual sola que detrás del proxy.} Cuatro payloads de
    Command Injection (\texttt{cmdi\_112}, \texttt{cmdi\_134}, \texttt{cmdi\_31} y
    \texttt{cmdi\_38}) los bloquea la Configuración~3b y no los bloquean ni la 2 ni la 3a. Con
    un servidor de captura y cinco payloads se comprobó que el cuerpo y el objetivo de la petición
    llegan idénticos por ambas rutas y que el proxy agrega cuatro cabeceras \texttt{X-*} y quita el
    puerto de \texttt{Host}; al
    enviar los cuatro payloads a la Configuración~2 con esas cabeceras, el paquete 0.1.1 los
    bloqueó en las tres repeticiones, y sin ellas no bloqueó ninguno, mientras que agregar además
    el \texttt{Host} sin puerto no cambió ningún veredicto. Con las cabeceras, la única
    característica del vector que cambia es el recuento de cabeceras fuera de la lista estándar
    (\texttt{unusual\_headers\_count}, de 0 a 4), que usa el Random Forest. La diferencia queda
    explicada por el proxy: la Configuración~2 y la librería dentro de la 3b no son la misma
    medición, y el veredicto de la librería depende en parte de las cabeceras que agrega el
    despliegue.
    \item \textbf{Causa del salto de Command Injection en la Configuración~3a.} Con el ejecutor
    anterior sobre la pila actual el WAF solo bloqueó 136 de 200 (68.0~\%), con 52 de 105 payloads
    de separadores ya codificados (49.5~\%); con el ejecutor corregido, 164 de 200 (82.0~\%) y 80 de
    105 (76.2~\%). La diferencia se atribuye a la recodificación de separadores ya codificados, y
    con el ejecutor corregido el WAF sigue bloqueando menos los codificados (76.2~\%) que los
    no codificados (88.4~\%, 84 de 95), una diferencia de 12.2 puntos porcentuales (prueba exacta de
    Fisher bilateral, $p = 0.028$) que no se atribuye a la codificación como causa, porque los dos
    subconjuntos no están controlados por otras características del payload.
        \item \textbf{Comparación con el WAF no neutral.} El corpus se construyó a partir de familias
    de payload que también están representadas en el entrenamiento de la librería, mientras que el
    WAF es un conjunto de reglas genérico en nivel de paranoia 1; la comparación entre ambos no es
    neutral y favorece a la librería por construcción. No se midieron otros niveles de paranoia ni
    un corpus ajeno al entrenamiento.
    \item \textbf{Repeticiones de un proceso determinista.} Las tres corridas por configuración
    (9{,}600 envíos de payload en total) no cambiaron el veredicto de ningún payload; confirman que
    el procedimiento es determinista y no estiman la variabilidad entre corpus ni entre
    construcciones del contenedor. La incertidumbre de muestreo la expresan los intervalos de
    Wilson.
    \item \textbf{Etiqueta flotante de la imagen del WAF.} \texttt{docker-compose.yml} referencia la
    imagen \texttt{owasp/modsecurity-crs:nginx-alpine} sin digest; el digest medido
    (\texttt{sha256:f12da37e...}) quedó registrado, pero la etiqueta resuelve hoy a un manifiesto
    distinto, de modo que repetir la medición sin fijar el digest puede cambiar el WAF.
    \item \textbf{La medición con el paquete 0.1.1 no es independiente.} Los defectos se
    encontraron analizando los fallos de este mismo corpus de la Ronda~5, y la medición que los da
    por corregidos (200 de 200 en Path Traversal) se hizo sobre ese corpus. Un resultado obtenido
    sobre los datos que motivaron la corrección es optimista por construcción; confirmarlo exige un
    corpus nuevo, generado con otra semilla y filtrado de la misma manera.
    \item \textbf{Saturación de casi-duplicados.} El corpus no se deduplicó más allá de los
    registros idénticos. La detección por MinHash, descrita en la sección VII.4 del trabajo de
    graduación de Diego Valenzuela \cite{valenzuela2026ia}, encontró saturación en las cinco clases: SQLi tiene
    aproximadamente 395 veces más pares casi duplicados que filas (dominada por la generación por
    plantillas de sqlmap y CAPEC) y la clase legítima, aproximadamente 9.5 veces. SQLi es el
    59.3~\% del conjunto de prueba. Con una partición aleatoria estratificada, variantes casi
    idénticas pueden quedar en entrenamiento y en prueba, lo que puede inflar las métricas, sobre
    todo las de SQLi; este trabajo no cuantificó ese efecto. El análisis por distancia de
    Levenshtein (211 pares, todos de la fuente CAPEC) lo subestimó, porque solo compara consultas de
    1 a 100 caracteres y hasta 10{,}000 comparaciones por clase.
    \item \textbf{Alcance de la evidencia sobre variantes desconocidas.} La detección de
    variantes no vistas de los vectores cubiertos se mide a nivel del sistema con el corpus limpio
    M2. El recall del Isolation Forest se mide sobre las mismas cuatro categorías que ve el Random
    Forest, por lo que no aísla la contribución del detector de anomalías a esa capacidad.
    \item \textbf{Validez externa del entorno simulado.} El entorno controlado no reproduce el
    tráfico de una startup en producción. La validación en aplicaciones de terceros reduce ese
    riesgo, pero no lo elimina.
\end{itemize}

% ------------------------------------------------------------------------------
\section{Restricciones y condiciones de ejecución}
\label{sec:met-restricciones}

El proyecto se ejecutó bajo las restricciones de tiempo (mayo a octubre de 2026), presupuesto y
recursos humanos descritas en el capítulo de Alcance. Además, las siguientes condiciones afectan
el alcance de la evidencia disponible para el capítulo de Resultados y se ejecutarán solo como
trabajo futuro:

\begin{itemize}
        \item \textbf{Explotabilidad de XSS confirmada en un solo caso.} La cobertura de XSS se midió
    sobre un corpus propio de 200 payloads, 64 de ellos dirigidos al campo vulnerable
    (\texttt{bio}); la explotabilidad confirmada sigue descansando en un único caso (ZAP,
    \texttt{POST /profile}, campo \texttt{bio}), que no se amplió. Que ZAP reporte un solo caso
    refleja la aplicación evaluada, no un límite de la herramienta.
    \item \textbf{Benchmark no repetido con el modelo sincronizado.} Ver la amenaza a la validez
    correspondiente en la Sección~\ref{sec:met-validez}.
    \item \textbf{Sin un escenario sostenido de condición única.} Cada repetición del benchmark
    dura 75~s. Se ejecutó además una corrida de monitoreo continuo de 41.9~minutos (374 muestras,
    con un pico de 494.8~MiB en el contenedor), pero esa serie mezcla condiciones distintas, y el
    indicador heurístico \texttt{possibleLeak}, que compara las dos mitades de la serie, no es
    concluyente por esa razón. No se ejecutó un escenario de 30~minutos o más con una sola condición
    estable, que es el que permitiría descartar una fuga de memoria.
    \item \textbf{Validación con startups no realizada.} No se concretó un acuerdo de
    colaboración en el período de pruebas; la validación en entornos productivos queda fuera de
    la evidencia del proyecto.
\end{itemize}

El Cuadro~\ref{tab:met-riesgos} actualiza los riesgos metodológicos del protocolo con su
estado al cierre de la ejecución.

\begin{longtable}{|L{3.6cm}|L{2.3cm}|L{6.9cm}|}
\hline
\textbf{Riesgo del protocolo} & \textbf{Estado} & \textbf{Mitigación aplicada} \\ \hline
\endfirsthead
\multicolumn{3}{l}{\footnotesize\textit{(continuación)}} \\ \hline
\textbf{Riesgo del protocolo} & \textbf{Estado} & \textbf{Mitigación aplicada} \\ \hline
\endhead
\multicolumn{3}{r}{\footnotesize\textit{Continúa en la página siguiente}} \\
\endfoot
\endlastfoot
Ningún modelo alcanza F1 $\geq$0.80 & No materializado & Datos reales adicionales en lugar de muestras sintéticas para Command Injection; lectura de prueba con verificación de hash \\ \hline
Fuga de datos entre particiones & Materializado en el corpus de evaluación, no en las particiones & Auditoría exacta y por plantilla; corpus limpio con filtro de exclusión \\ \hline
Diferencia de paridad ONNX $>$0.1~\% & No materializado & Prueba de paridad obligatoria en CI \\ \hline
Falsos positivos ante tráfico real & Materializado & El Isolation Forest no bloquea; estudio con navegadores; modo \texttt{monitor} documentado para una primera instalación \\ \hline
Imposibilidad de acuerdo con startups & Materializado & Sustitución por cuatro aplicaciones de terceros \\ \hline
Extracción de características $>$5~ms & No materializado & Extractor único en TypeScript con benchmark propio \\ \hline
\caption[Estado de los riesgos metodológicos]{Estado de los riesgos metodológicos del protocolo al cierre de la ejecución}
\label{tab:met-riesgos}
\end{longtable}
{\footnotesize Nota. Elaboración propia.}

La viabilidad del método descansa en que todos los procedimientos usan herramientas gratuitas
y scripts versionados, se ejecutan en hardware local o en ejecutores públicos de integración
continua y dejan sus datos crudos en los repositorios del proyecto. Esto permite que otro
investigador los repita sin depender del entorno de los autores.

\ifdefined\modoensamble
\else
	\section*{Referencias}
	\printbibliography[heading=none]
	\end{document}
\fi
